% concatenated from GaVe_Revision_final.tex
\documentclass[11pt]{amsart}
% \usepackage{times} %mette times new roman come font

\usepackage{amsfonts, enumerate} %amsfonts %queste annotazioni mi dicono il parametro che potrei usare, ad esempio in questo caso equivale a dire che potrei togliere "enumerate"
\usepackage{amssymb}
\usepackage[final]{graphicx, graphics,stmaryrd} %graphicx

%\input xy %non ho ancora scoperto cos'è ma credo serva
%\xyoption{all} %come sopra
%\usepackage[all]{xy}

\usepackage[margin=3cm]{geometry} 
%\usepackage{amscd}


%\usepackage[breaklinks,bookmarksopen,bookmarksnumbered]{hyperref}

%\usepackage{fullpage}
%\usepackage{tensor}
\usepackage[english]{babel}
\usepackage[utf8]{inputenc}
\usepackage{amsthm}
\usepackage{graphics}
\usepackage{amsmath}
\numberwithin{equation}{section} %numerazione delle equazioni 
\usepackage{amstext}

\usepackage[safe,extra]{tipa}
%\usepackage{showkeys}  %mostra o nasconde i nomi dei label sul file di output
\usepackage{multirow}
%\usepackage{enumitem}
%\usepackage{refcheck}

\usepackage[colorlinks=true,urlcolor=blue,
citecolor=red,linkcolor=blue,linktocpage,pdfpagelabels,
bookmarksnumbered,bookmarksopen]{hyperref}
\usepackage[hyperpageref]{backref}


%\usepackage{relsize} % per fare scritte matematica più grosse col comando \mathlarger dà grossi problemi con hyperref
\usepackage{bm} %simboli BOLD in math mode
\usepackage{stackrel} %per scrivere cose impilate in math mode
\usepackage{tikz-cd} %diagrammi
\usepackage{mathtools} %%???

\newtheorem{teo}{Theorem}[section]
\newtheorem{coro}[teo]{Corollary}
\newtheorem{prop}[teo]{Proposition}
\newtheorem{lemma}[teo]{Lemma}
\newtheorem*{props}{Proposition}
\newtheorem*{teon}{Theorem}
%\newtheorem*{conj}{Conjecture}
%\newtheorem*{que}{Question}

\theoremstyle{definition}
\newtheorem{defin}[teo]{Definition}
\newtheorem{rmk}[teo]{Remark}
\newtheorem{ex}[teo]{Example}



\newcommand{\mb}{\mathbb}
\newcommand{\mc}{\mathcal}
\newcommand{\msf}{\mathsf}
\newcommand{\mfk}{\mathfrak}

\def\Oo{\mathcal O}
\def\R{\mathbb{R}}
\def\C{\mathbb C}
\def\N{\mathbb N}
\def\Z{\mathbb Z}
\def\K{\mathbb K}
\def\Q{\mathbb Q}
\def\p{\mathbb P}
\def\cc{\mathcal C}
\def\E{\mc E}
\def\M{\mathcal{M}}
\def\A{\mathcal A}
\def\pp{\mathfrak P}
\def\mm{\bm{\mathsf{M}}}
\def\rr{\bm{\mathsf{R}}}
\def\U{\mathcal U}
\def\cc{\mathcal{C}}
\def\D{\mathcal{D}}
\def\s{\mc{S}}
\def\H{\mb{H}}
\def\J{\mc{J}}
\def\I{\mc{I}}
\def\F{\mfk{F}}
\def\g{\mfk{g}}
\def\d{\mathrm{d}}
\def\aa{\mc{A}}
\def\nn{\mc{N}}
\newcommand{\al}{\alpha}
\newcommand{\be}{\beta}
\def\G{\mathbb{G}}



\newcommand{\m}{\mbox}
\newcommand{\gra}{\textbf}
\newcommand{\cor}{\textit}
\newcommand{\sott}{\underline}

\newcommand{\fine}{\qed\newline}
\newcommand{\xx}{\otimes}


\newcommand{\longhookrightarrow}{\ensuremath{\lhook\joinrel\relbar\joinrel\rightarrow}}

\newcommand{\GG}{\mb G}
\newcommand{\we}{\wedge}
\newcommand{\Lim}{\varprojlim}
\newcommand{\Colim}{\varinjlim}



\DeclareMathOperator{\age}{age}
\DeclareMathOperator{\fix}{Fix}
\DeclareMathOperator{\defo}{Def}
\DeclareMathOperator{\Aut}{Aut}
\DeclareMathOperator{\diag}{Diag}
\DeclareMathOperator{\ord}{ord}
\DeclareMathOperator{\im}{Im}
\DeclareMathOperator{\id}{id}
\DeclareMathOperator{\sing}{Sing}
\DeclareMathOperator{\tr}{Tr}
\DeclareMathOperator{\sym}{Sym}
\DeclareMathOperator{\End}{End}
\DeclareMathOperator{\Ho}{Hom}
\DeclareMathOperator{\Ker}{Ker}
\DeclareMathOperator{\Spec}{Spec}
\DeclareMathOperator{\Proj}{Proj}
\DeclareMathOperator{\Etspec}{EtSpec}
\DeclareMathOperator{\ext}{Ext}
\DeclareMathOperator{\GL}{GL}
\DeclareMathOperator{\NE}{ne}
\DeclareMathOperator{\QR}{QR}
\DeclareMathOperator{\Pic}{Pic}
\DeclareMathOperator{\cut}{cut}
\DeclareMathOperator{\Rep}{Rep}
\DeclareMathOperator{\Div}{Div}
\DeclareMathOperator{\codim}{Codim}
\DeclareMathOperator{\rk}{rk}
\DeclareMathOperator{\nc}{nc}
\DeclareMathOperator{\GCD}{GCD}
\DeclareMathOperator{\T}{T}
\DeclareMathOperator{\Bl}{Bl}
\DeclareMathOperator{\surj}{Surj}
\DeclareMathOperator{\Ob}{Ob}
\DeclareMathOperator{\Mor}{Mor}
\DeclareMathOperator{\dom}{dom}
\DeclareMathOperator{\cod}{cod}
\DeclareMathOperator{\op}{op}
\DeclareMathOperator{\Char}{char}
\DeclareMathOperator{\Flat}{Flat}
\DeclareMathOperator{\Filt}{Filt}
\DeclareMathOperator{\Coflat}{CoFlat}
\DeclareMathOperator{\uspec}{\underline{Spec}}
\DeclareMathOperator{\uproj}{\underline{Proj}}
\DeclareMathOperator{\vol}{Vol}
\DeclareMathOperator{\inv}{inv}
\DeclareMathOperator{\Bun}{Bun}
\DeclareMathOperator{\Root}{Root}
\DeclareMathOperator{\gen}{gen}
\DeclareMathOperator{\lcm}{lcm}
\DeclareMathOperator{\Def}{Def}
\DeclareMathOperator{\Fun}{Fun}
\DeclareMathOperator{\Tor}{Tor}
\DeclareMathOperator{\Sym}{Sym}
\DeclareMathOperator{\ev}{ev}
\DeclareMathOperator{\inn}{Inn}
\DeclareMathOperator{\out}{Out}
\DeclareMathOperator{\sub}{Sub}
\DeclareMathOperator{\lex}{Lex}
\DeclareMathOperator{\colex}{CoLex}
\DeclareMathOperator{\Mod}{Mod}
\DeclareMathOperator{\Nat}{Nat}
\DeclareMathOperator{\Cl}{Cl}
\DeclareMathOperator{\Zar}{Zar}
\DeclareMathOperator{\Lie}{Lie}
\DeclareMathOperator{\supp}{supp}
\DeclareMathOperator{\rat}{rat}
\DeclareMathOperator{\et}{et}
\def\h{\mfk{h}}
\newcommand{\de}{\partial}
\DeclareMathOperator{\Lin}{Lin}
\DeclareMathOperator{\loc}{loc}
\DeclareMathOperator{\Exp}{Exp}
\DeclareMathOperator{\El}{El}
\DeclareMathOperator{\Geod}{Geod}
\DeclareMathOperator{\Adm}{Adm}
\DeclareMathOperator{\Opt}{Opt}
\DeclareMathOperator{\Log}{Log}
\DeclareMathOperator{\ce}{ce}
\DeclareMathOperator{\se}{se}
\DeclareMathOperator{\In}{in}
\DeclareMathOperator{\Out}{out}

\def\din{\Omega^{\In}}
\def\dout{\Omega^{\Out}}




%\input{newcommands_eng.tex}
%\title{titolo}

\makeatletter
\@namedef{subjclassname@2020}{\textup{2020} Mathematics Subject Classification}
\makeatother
%\usepackage{refcheck}

 \begin{document}
 	\title[Crtitical Concave-convex problems in Carnot groups]{Critical concave-convex problems in Carnot groups}
 	
 
 	\author[M.\,Galeotti]{Mattia Galeotti}
 	\author[E.\,Vecchi]{Eugenio Vecchi}
 	
  	
 	\address[M.\,Galeotti]{Dipartimento di Matematica
 		\newline\indent Università  di Bologna \newline\indent
 		Piazza di Porta San Donato 5, 40126 Bologna, Italy}
 	\email{mattia.galeotti4@unibo.it}
 	
 	\address[E.\,Vecchi]{Dipartimento di Matematica
 		\newline\indent Università  di Bologna \newline\indent
 		Piazza di Porta San Donato 5, 40126 Bologna, Italy}
 	\email{eugenio.vecchi2@unibo.it}
 	
 	\keywords{PDEs on Carnot groups, Critical equations, Sublinear equations}
 	
 	\subjclass[2020]{35R03, 35B33, 35H20, 35J70}
 	%35R03 PDEs on Heisenberg groups, Lie groups, Carnot groups, etc.
 	%35J75 Singular elliptic equations
 	%35J20 Variational methods for second-order elliptic equations
 	%35B33 Critical exponents in context of PDEs
 	%35J70 Degenerate elliptic equations
 	%35H20 Subelliptic equations
  	%35B09 Positive solutions to PDEs
 	
 	\date{\today}
 	
 	\thanks{E.V. and M.G.are
 		members of GNAMPA-Indam. E.V. is
 		partially 
 		supported by the PRIN 2022 project 2022R537CS \emph{$NO^3$ - Nodal Optimization, NOnlinear elliptic equations, NOnlocal geometric problems, with a focus on regularity}, founded by the European Union - Next Generation EU and by the Indam-GNAMPA project CUP E5324001950001 - {\em Problemi singolari e degeneri: esistenza, unicità e analisi delle proprietà qualitative delle soluzioni}.
 		M.G. is supported by the National Recovery and Resilience Plan (NRRP), research funded by the European Union – NextGenerationEU.
 		CUP D93C22000930002, {\em A multiscale integrated approach to the study of the nervous system in health and disease (MNESYS)}}
 	
 	
 
 \begin{abstract}
 We consider a model Dirichlet problem with concave-convex and critical nonlinearity settled in Carnot groups. Our aim is to prove the existence of two positve solutions in the spirit of a famous result by Ambrosetti, Brezis and Cerami. To this aim we use a variational Perron method combined with proper estimates of a family of functions which are minimizers of the relevant Sobolev inequality. Due to the lack of boundary regularity, we also have to be careful while proving that the first solution found is a local minimizer in the proper topology.
 \end{abstract}	
 
 	
\maketitle

\section{Introduction}
Let $\mathbb{G}$ be a Carnot group and let $\Omega \subset \mathbb{G}$ be an open, bounded and connected set with Lipschitz boundary $\partial \Omega$. Let $q \in (0,1)$, let $2^{\star}_{Q}:=\tfrac{2Q}{Q-2}$ be the critical Sobolev exponent related to the Sobolev inequality in $\mathbb{G}$, and let $\lambda >0$. We consider the following Dirichlet boundary value problem
\begin{equation}\tag{{$\mathrm{P}_{\lambda}$}}\label{EqProblem}
	\left\{\begin{array}{rll}
		-\Delta_{\mathbb{G}}u& = \lambda\, u^{q} + u^{2^{\star}_{Q}-1} & \textrm{ in } \Omega,\\
		u&>0 & \textrm{ in } \Omega,\\
		u&=0 & \textrm{ on } \partial \Omega.
	\end{array}\right.
\end{equation}

We stress that $-\Delta_{\mathbb{G}}$ denotes here the sub-laplacian on $\mathbb{G}$ which is a second-order differential operator with non-negative characteristic form that can be explicitly expressed as a sum of squares of vector fields satisfying the H\"{o}rmander condition, see e.g. \cite{Hormander}. We refer to Section~\ref{sec:Prel} for more details, including the Folland-Stein Sobolev spaces we will work with.\\
%
%Along the paper it will sometimes be useful to denote the above problem as \eqref{EqProblem}$_\lambda$ 
%to make it clear the choice of the parameter.
We immediately state the main result of this paper. In what follows,
we refer to Definition \ref{def:weak_sub_super_sol} for the precise definition of \emph{weak solution}
of \eqref{EqProblem}.
\begin{teo} \label{thm:main}
	Let $\Omega\subset\mathbb{G}$ be an open, bounded and connected
	set
	with Lipschitz boundary $\partial \Omega$, and let $q\in (0,1)$. Then, there exists $\Lambda > 0$
	such that
	\begin{itemize}
		\item[A)]problem \eqref{EqProblem} does not admit weak solutions
		for every $\lambda>\Lambda$;
		\item[B)] problem \eqref{EqProblem} admits at least one weak solution for every $\lambda \in (0,\Lambda]$;
		\item[C)] problem \eqref{EqProblem} admits at least two weak solutions for every $0<\lambda<\Lambda$.
	\end{itemize}
\end{teo}

The above theorem is the natural generalization to Carnot groups of classical results of \cite{ABC}. We refer e.g. to \cite{BESS,CCP,CoPe,BV} for further generalizations. 


The interest in studying existence of positive solutions to critical problems in the Carnot group setting,
 is in the geometric significance of the purely critical problem in the model case of the Heisenberg group.
 Indeed, when $\lambda =0$ and $\Omega = \mathbb{H}^n$, the problem \eqref{EqProblem} becomes the famous CR-Yamabe problem studied by Jerison and Lee \cite{JerisonLee,JerisonLee2,JerisonLee3}. The problem  we are interested in is settled on bounded domains, where tipycally one can prove non-existence of positive solutions, at least in star-shaped domains, by appealing suitable versions of the Pohozaev identity. Because of this, the seminal paper by Brezis and Nirenberg \cite{BN} showed that adding a perturbative term, linear in \cite{BN}, but subsequently extended to much more general perturbations, may allow to prove the existence of one or more positive solutions. A crucial tool in the argument performed in \cite{BN} is provided by the use of the Aubin-Talenti functions, whose analogue in $\mathbb{H}^n$ made its appearance in \cite{JerisonLee2}. This was a key ingredient which gave rise to a prolific study of critical problems in $\mathbb{H}^n$, see e.g. \cite{Citti, GaLa, Ugu1, CitUg, MaUg, FelliUgu, MaMaPi, PaPiTe}.

As long as one needs explicit knowledge of proper replacements of the Aubin-Talenti functions, the only other sub-Riemannian structure where they are known is that of groups of Iwasawa type, see \cite{GaroVa2, GaroVa}. As far as we know, there are no other structures, nor Sobolev inequalities with $p\neq 2$, for which the minimizers are explicitly known. 
On the other hand, since the best constant in the Sobolev inequality is achieved in all Carnot groups (see \cite{GaroVa2}), it has been proved to be enough to know  the asymptotic behaviour at infinity of the minimizers. This is now known for $p\neq 2$ as well, see \cite{Loiudice3}, and it paved the way for a series of 
existence, multiplicity or non-existence of positive solutions for critical problems  à la Br\'{e}zis-Nirenberg in $\mathbb{G}$: we refer e.g. to \cite{BoUg, Loiudice1, Loiudice2, Loiudice4, BiGaVe}.

\medskip

Let us now briefly describe the proof of Theorem \ref{thm:main}: 
%the argument performed in the Euclidean case by , suitably adapted to the Carnot group setting, and it consists of several technical steps that we list here below:
\begin{itemize}
	\item in Theorem \ref{thm:Existence_First} we prove the existence of a first solution by means of a variational Perron method which transfers the approach of Struwe \cite{Struwe} to the Carnot group setting. In particular, setting
	\begin{equation*}
	\Lambda := \sup \{ \lambda >0: \eqref{EqProblem} \textrm{ admits a weak solution}\},
\end{equation*}
	\noindent we show first that $0<\Lambda<+\infty$, and this immediately provides a threshold for the non-existence of weak solutions. Once this is done, we use the unique solution of the purely sublinear problem \eqref{eq:Sublinear_Problem} as a weak subsolution and we construct a weak supersolution for fixed $\lambda$ by using the weak solution for a bigger $\lambda'$; 
	\item we show that for $\lambda \in (0,\Lambda)$ the first solution obtained as described before is a local minimizer in the natural topology associated with problem \eqref{EqProblem}, see Lemma \ref{lem:locmin}. We stress here that in \cite{ABC} the authors made use of a famous result by Brezis and Nirenberg \cite{BNH1C1} which does not have an analog in the Carnot group setting. This is due to the fact that $C^{1,\alpha}$ regularity up to the boundary is still a delicate issue at the so called characteristic points: the first obstructions have been observed by Jerison \cite{Jerison,Jerison2}, but this is still an active field of research, see e.g. \cite{BaCiCu,BaGaMu,AbTr}. For this reason we follow here a more variational approach based on a paper by Alama \cite{Alama}, already used in a different setting in \cite{AbDiVa};
	\item we prove the existence of a second solution following an argument originally due to Tarantello \cite{Tarantello}: this combines the Ekeland variational principle \cite{Ekeland} with the fine asymptotic expansions proved in \cite{Loiudice1}.
\end{itemize} 

\medskip

We stress that the multiplicity result obtained in Theorem \ref{thm:main} can be easily extended to cover the convex-case of a Sobolev sub-critical nonlinearity.

\medskip

The paper is organized as follows: in Section \ref{sec:Prel} we recall the basic facts on Carnot groups and we set the variational functional setting necessary for the study of \eqref{EqProblem}. We also recall the basic result regarding the purely sublinear problems, like existence and uniqueness of a positive solution and a comparison principle resembling the classical one. In Section \ref{sec:First_Solution} we prove the existence of a first solution as described before, while the existence of a second solution (for $\lambda \in (0,\Lambda)$) is postponed to the final Section \ref{sec:Second_Solution}.





\section{Preliminaries}\label{sec:Prel}
In this section we collect all the relevant notations, definitions and preliminaries needed in the rest of the paper.
\subsection{Carnot groups}
A Carnot group $\mathbb{G}=(\mathbb{R}^{N},\diamond)$ of step $k$ is a connected, simply connected Lie group whose finite dimensional Lie algebra $\mathfrak{g}$ of left-invariant (w.r.t. $\diamond$) vector fields admits a stratification of step $k$, namely there exist $k$ linear subspaces $\mathfrak{g}_1, \ldots, \mathfrak{g}_k$ such that
\begin{equation*}
	\mathfrak{g} = \mathfrak{g}_{1}\oplus \ldots \oplus \mathfrak{g}_{k}, \qquad [\mathfrak{g}_1,\mathfrak{g}_i]=\mathfrak{g}_{i+1}, \qquad \mathfrak{g}_{k}\neq \{0\}, \qquad \mathfrak{g}_i = \{0\} \textrm{ for all } i>k.
\end{equation*}
In particular, this implies that Carnot groups are a special instance of graded groups.\\
We call $\mathfrak{g}_1$ the horizontal layer. We denote by $X_1, \ldots, X_{N}$ a basis of left-invariant vector fields of $\mathfrak{g}$ such that the following holds:
\begin{itemize}
	\item $X_1, \ldots, X_{m_1}$ is an orthonormal basis of $\mathfrak{g}_{1}$ w.r.t. the scalar product $\langle \cdot, \cdot \rangle_{\mathfrak{g}_{1}}$;
	\item for every $1<i\leq k$, $X_{m_{i-1}+1},\ldots, X_{m_{i}}$ is a basis of of $\mathfrak{g}_{i}$;
	\item $m_{0}=0$ and $n_i := m_{i}-m_{i-1} = \dim \mathfrak{g}_i$ for every $1\leq i\leq k$;
	\item $m_1 + \ldots + m_k = N$.
\end{itemize}
We define the homogeneous dimension of $\mathbb{G}$ as
\begin{equation*}%\label{eq:Def_Q}
	Q:= \sum_{i=1}^{k}i \cdot n_{i}.
\end{equation*}
%We notice that $N$ is the topological dimension of $\mathbb{G}$, but we can also define its {\em homogeneous dimension} $Q$ as follows
%\begin{equation}\label{eq:Def_Q}
%	Q:= \sum_{i=1}^{k}i \, n_{i}.
%\end{equation}
%We notice that $N \leq Q$ and that $N=Q$ if and only if $\mathbb{G}$ is the classical Euclidean group $(\mathbb{R}^{N},+)$. In particular, this is the only possible case whenever $Q \leq 3$.\\
%Since the exponential map is a one-to-one diffeomorphism from $\mathfrak{g}$ to $\mathbb{G}$, any point $g\in \mathbb{G}$ can be uniquely written in exponential coordinates as
%\begin{equation*}
%	g = g_1 X_1 + \ldots + g_N X_N = (g_1, \ldots, g_N).
%\end{equation*}
The left translations $\tau : \mathbb{G}\to \mathbb{G}$ provides a family of automorphisms of $\mathbb{G}$, and are defined as follows
\begin{equation}\label{eq:left-translation}
	\tau_{h}(g):= h \diamond g, \quad \textrm{ for a given $h\in \mathbb{G}$,} 
\end{equation}
The anisotropic dilations $\delta_{\lambda}: \mathbb{G}\to \mathbb{G}$ of $\mathbb{G}$ are instead defined as 
\begin{equation}\label{eq:dilation}
	\delta_{\lambda}(g) = \left(\lambda^{\alpha_1}g_1, \ldots, \lambda^{\alpha_{N}}g_N \right), \quad \textrm{ for every } \lambda >0,
\end{equation}
\noindent where $\alpha_{j} = i$ if $m_{i-1}<j\leq m_{i}$.
We notice that $Q = \alpha_1 + \ldots + \alpha_{N}$.\\
% The null element of $\diamond$ is the identity $0 = (0,\ldots,0)$ and the inverse of a certain $g \neq 0$ is usually denoted by $g^{-1}$.\\
%The group operation $\diamond$ can also be used to define a further family of automorphisms of $\mathbb{G}$ known as {\em left translations} $\tau : \mathbb{G}\to \mathbb{G}$. More precisely, given a base point $h\in \mathbb{G}$, we define 
%\begin{equation}\label{eq:left-translation}
%	\tau_{h}(g):= h \diamond g.
%\end{equation}

\medskip

%It is possible to endow a Carnot group $\mathbb{G}$ with a richer structure. Firstly, it is possible to define a scalar product $\langle \cdot, \cdot \rangle_{\mathfrak{g}_{1}}$ such that the basis $\{X_1, \ldots, X_{m_1}\}$ of the horizontal layer becomes orthonormal, and this provides a sub-Riemannian structure over $\mathbb{G}$. We notice that one can also provide a purely Riemannian structure defining a scalar product on the Lie algebra $\mathfrak{g}$ making the entire basis $\{X_1,\ldots,X_N\}$ orthonormal.
%The sub-Riemannian Carnot group $\mathbb{G}$ can also be endowed with an intrinsic metric structure by means of the so called {\em Carnot-Carath\'{e}odory} ($CC$ in short) distance. It is well known that with this distance, these spaces are not Riemannian at any scale, see e.g. \cite{Semmes}.


Given a smooth horizontal vector field $V=v_1 X_1 + \ldots + v_{m_1}X_{m_1}$, we define its horizontal divergence as
\begin{equation}\label{eq:horizontal_divergence}
	\mathrm{div}_{\mathbb{G}}V := X_{1}v_1 + \ldots + X_{m_1}v_{m_1}.
\end{equation}
Moreover, given a smooth enough scalar-valued function $u:\mathbb{G}\to \mathbb{R}$, we can define
the horizontal gradient of $u$ as
\begin{equation}\label{eq:Horizontal_grad}
	\nabla_{\mathbb{G}}u := (X_{1}u, \ldots, X_{m_1}u),
\end{equation}
\noindent and the sub-Laplacian of $u$ as
\begin{equation}\label{eq:horizontal_lap}
	\Delta_{\mathbb{G}}u := \mathrm{div}_{\mathbb{G}}(\nabla_{\mathbb{G}}u) = X_1^2 u + \ldots + X_{m_1}^{2} u,
\end{equation}
%We point out that both the divergence $\mathrm{div}_{\mathbb{G}}$ and the horizontal gradient $\nabla_{\mathbb{G}}$ (and \cor{a fortiori} the sub-Laplacian $\Delta_{\mathbb{G}}$)
%are independent of the choice of the base on $\g_1$: in fact, they are intrisically
%associated to the sub-Riemannian structure on $\mathbb{G}$.

%We notice that both the horizontal gradient $\nabla_{\mathbb{G}}$ and the sub-Laplacian $\Delta_{\mathbb{G}}$ are left-invariant operators, i.e.
%\begin{equation}
%	\nabla_{\mathbb{G}}(u\circ \tau_{h}) = (\nabla_{\mathbb{G}}u)\circ \tau_{h} \quad \textrm{ and } \quad \Delta_{\mathbb{G}}(u \circ \tau_{h}) = (\Delta_{\mathbb{G}}u)\circ \tau_{h}, \quad \textrm{ for every } h \in \mathbb{G},
%\end{equation}
%\noindent and they are, respectively, homogeneous of degree one and two w.r.t. the family of dilations $\delta_{\lambda}$ defined in \eqref{eq:dilation}, namely
%\begin{equation}
%	\nabla_{\mathbb{G}}(u\circ \delta_{\lambda})= \lambda (\nabla_{\mathbb{G}}u)\circ \delta_{\lambda} \quad \textrm{ and } \quad \lambda^{2}(\Delta_{\mathbb{G}}u)\circ \delta_{\lambda}, \quad \textrm{ for every } \lambda >0.
%\end{equation}
The Lebesgue measure $\mathcal{L}^{N}$ coincides with the Haar measure of $\mathbb{G}$ and hence is left-invariant and satisfies the following scaling property:
\begin{equation}
	\mathcal{L}^{N}(\delta_{\lambda}(E)) = \lambda^{Q} \mathcal{L}^{N}(E) \quad \textrm{ for every measurable set } E \subset \mathbb{G}.	
\end{equation}
Every integral in this manuscript has to be understood
with respect to the Haar measure, unless otherwise stated.

%Moreover, the homogeneity of the $X_i$'s implies that the (formal)
%adjoint of $X_i$ in the Lebesgue space $L^2(\mathbb{G})$ is precisely $-X_i$ (for $i\leq i\leq m_1$), that is,
%\begin{equation} \label{eq:XiSA}
% \int_{\mathbb{G}}(X_i\varphi)\psi = -\int_{\mathbb{G}}\varphi(X_i\psi)\quad\text{for
% every $1\leq i\leq m_1$}.
%\end{equation}
%In particular, $-\Delta_\mathbb{G}$ is a \emph{self-adjoint operator}.
\medskip

Every Carnot group can be endowed with several homogeneous norms. A homogeneous (quasi)norm $\rho : \mathbb{G} \to \mathbb{R}$ is a non-negative function further satisfying the following properties:
\begin{itemize}
	\item $\rho(g)=0$ if and only if $g=0$;
	\item $\rho(\delta_{\lambda}(g)) = \lambda \, \rho(g)$ for every $g \in \mathbb{G}$ and for every $\lambda >0$;
	\item $\rho(h\diamond g) \leq C \left(\rho(h) + \rho(g)\right)$ for every $g,h \in \mathbb{G}$ and for some constant $C \geq 1$.
\end{itemize}
By a famous result of Folland \cite{Folland}, there exists a homogeneous norm $|\cdot|_{\mathbb{G}}$ on $\mathbb{G}$ and a positive constant $C_Q>0$, depending only on $Q$, such that the function
\begin{equation}\label{eq:def_Gamma}
	\Gamma_{h}(g):= \dfrac{C_Q}{|h^{-1}\diamond g|^{Q-2}_{\mathbb{G}}}, \quad \textrm{ with } Q \geq 3,
\end{equation}
\noindent is a fundamental solution of $-\Delta_{\mathbb{G}}$ with pole at $h \in \mathbb{G}$.
Moreover, homogeneous norms can be used to define distances as follows: 
\begin{equation*}
	d_{\rho}(g,h):= \rho(h^{-1}\diamond g).
\end{equation*}
We stress that there are other possible choices of distances (i.e. the so called $CC$-distance and many others) which are all equivalent.
%
%
%In any case, all these norms (and the relative distances) are equivalent and they all induce on $\mathbb{G}$ the Euclidean topology. For our purposes, we prefer to work with the homogeneous norm $|\cdot|_{\mathbb{G}}$ (and the associated distance $d_{\mathbb{G}}$) which provides the fundamental solution defined in \eqref{eq:def_Gamma}. 
Finally, we will denote by 
\begin{equation*}
	B_{r}(g_0) := \{ g \in \mathbb{G}: d_{\mathbb{G}}(g,g_0) = |g_{0}^{-1}\diamond g|_{\mathbb{G}} <r\},
\end{equation*} 
\noindent the open ball of radius $r>0$ and center $g_{0}\in \mathbb{G}$.\\
We refer e.g. to \cite{BLU} for a comprehensive introduction to the subject.

\medskip

\subsection{The functional setting}\label{sec:functional_setting}
Let $\mathcal{O} \subseteq \mathbb{G}$ be an open set.
For every $f \in C^{\infty}_{0}(\mathcal{O})$ there exists a positive constant $C_Q>0$ depending only on the homogeneous dimension $Q$ such that the following Sobolev inequality holds true
\begin{equation}\label{eq:Sobolev_Ineq}
	\|f\|^{2}_{L^{2_{Q}^{\star}}(\mathcal{O})} \leq C_{Q} \, \| |\nabla_{\mathbb{G}}f| \|^{2}_{L^{2}(\mathcal{O})},
\end{equation}
\noindent where 
\begin{equation}
	2_{Q}^{\star} := \dfrac{2Q}{Q-2},
\end{equation}
\noindent denotes the (sub-elliptic) critical Sobolev exponent.
Thanks to \eqref{eq:Sobolev_Ineq}, 
$\| |\nabla_{\mathbb{G}}f| \|_{L^{2}(\Omega)}$ provides a norm on the space $C^{\infty}_{0}(\Omega)$.
We define the Folland-Stein space $S^{1}_{0}(\mathcal{O})$ as the completion of $C^{\infty}_{0}(\mathcal{O})$ w.r.t.\,the above norm, and we set
$$\|u\|_{S^{1}_{0}(\mathcal{O})} = \||\nabla_\mathbb{G}u|\|_{L^2(\mathcal{O})}
\quad\text{for every $u\in S^{1}_{0}(\mathcal{O})$}.$$
We explicitly observe that, owing to \eqref{eq:Sobolev_Ineq}, we have
\begin{equation} \label{eq:explicitS01def}
 S_0^1(\mathcal{O}) = \big\{u\in L^{2_{Q}^{\star}}(\mathcal{O}):\,\text{$X_iu\in L^2(\mathcal{O})$
 for all $1\leq i\leq m_1$}\big\},
\end{equation}
where $X_1u,\ldots,X_{m_1}u$ are meant in the sense of distributions.\\
We now remind a couple of basic properties of $S_0^1(\mathcal{O})$ when $\mathcal{O}$ is an open and bounded set, see e.g. to \cite{FollandStein}.
\begin{itemize}
 \item $S_0^1(\mathcal{O})$ is endowed with a structure of real Hilbert space by the inner product
 $$\langle u,v\rangle_{S_0^{1}(\mathcal{O})} = \int_{\mathcal{O}}\langle \nabla_\mathbb{G}u,
 \nabla_\mathbb{G}v\rangle_{\mathfrak{g}_{1}}\qquad (u,v\in S_0^1(\mathcal{O})),$$
whose associated norm is precisely $\|\cdot\|_{S_0^1(\mathcal{O})}$.

 \item $S_0^1(\mathcal{O})$ is continuously embedded into $L^p(\mathcal{O})$ for every $1\leq p\leq 2_Q^\star$.
 Furthermore, this embedding turns out to be \emph{compact} when $1\leq p < 2_Q^\star$. 
% \item If $u\in S_0^1(\mathcal{O})$, then $u_+ = \max\{u,0\},\,u_- = \max\{-u,0\}\in S_0^1(\mathcal{O})$, and
% $$\nabla_\mathbb{G}u_+ = \nabla_\mathbb{G}u\cdot\chi_{\{u > 0\}}
% \quad \textrm{ and } \quad  \nabla_\mathbb{G}u_- = \nabla_\mathbb{G}u\cdot\chi_{\{u < 0\}}.
% $$
% In particular, for every $c \in\mathbb{R}$, we derive
% \begin{equation}\label{eq:Grad_and_level_set}
% 	\nabla_\mathbb{G}u = 0  \quad \textrm{a.e. on any level set } \{u = c\}.
% \end{equation}
 \end{itemize}
We refer e.g. to \cite{FollandStein} for more details.
%\begin{rmk} \label{rem:ConvergenceS01}
%On account of the above properties of $S_0^1(\Omega)$, it is possible to prove
%the following \emph{convergence result}, which will be repeatedly used in the sequel.
%
%Assume that $\{u_k\}_k\subseteq S_0^1(\Omega)$ is a \emph{bounded sequence}. Since $S_0^1(\Omega)$ is a real \emph{Hilbert space}, and since we have already pointed out that the embedding 
%  $$S_0^1(\Omega)\hookrightarrow L^p(\Omega)$$
%  is compact for every $1\leq p< 2^\star_Q$, as in the classical case we deduce that there exists a function $u\in S_0^1(\Omega)$ such that (up to a subsequence, and as $k\to+\infty$)
%  \vspace{0.1cm}
%  
%  \begin{itemize}
%  	\item  $u_k\to u$ \emph{weakly in $S_0^1(\Omega)$};
%  	\item $u_k\to u$ \emph{strongly} in $L^p(\Omega)$ for every $1\leq p< 2^\star_Q$;
%  	\item $u_k\to u$ \emph{pointwise a.e.\,in $\Omega$}.
%  \end{itemize}
%  On the other hand, since the embedding $S_0^1(\Omega)\hookrightarrow L^{2^*Q}(\Omega)$ is \emph{continuous}, we have that $u\in L^{2^\star_Q}(\Omega)$, and  $\{u_k\}_k$ is \emph{bounded also in $L^{2^\star_Q}(\Omega)$}; in particular, setting 
%  $$p_0 = \frac{2^\star_Q}{2^\star_Q-1}\in (1,2^\star_Q),$$
%  we have that $\{v_k = u_k^{2^\star_Q-1}\}_k$ is bounded in $L^{p_0}(\Omega)$. As a consequence of this fact, and since we know that $u_k\to u$ po\-int\-wise a.e.\,in $\Omega$, we deduce that $\{v_k\}_k$ \emph{converges weakly to $u^{2^\star_Q-1}$}
%  (as $k\to +\infty)$, and up to a subsequence) in $L^{p_0}(\Omega)$, that is, 
%  $$\int_{\Omega}u_k^{2^\star_Q-1}\varphi\,dx \to \int_{\Omega}u^{2^\star_Q-1}\varphi\,dx\quad\text{for every $\varphi\in L^{p_0'}(\Omega)
%  = L^{2^\star_Q}(\Omega)\supseteq S_0^1(\Omega)$}.$$
%\end{rmk}

\medskip

We are now ready to properly set the definition of weak sub/supersolution of \eqref{EqProblem}.
%However, we prefer to provide such definitions for the following slightly more general singular Dirichlet problem:
%\begin{equation}\label{eq:General_Singular_Problem}
%	\left\{\begin{array}{rl}
%		-\Delta_{\mathbb{G}}u = \lambda \, u^{q} + f(x,u)& \textrm{ in } \Omega,\\
%		u>0 & \textrm{ in } \Omega,\\
%		u=0 & \textrm{ on } \partial \Omega,
%	\end{array}\right.
%\end{equation}
%\noindent where $q \in (0,1)$, $\lambda >0$ and $f: \Omega \times (0,+\infty)\to \mathbb{R}$ is a Carath\'{e}odory function satisfying the following critical growth condition: there exists a positive constant $K_{f}>0$ such that
%\begin{equation}\label{eq:Growth_of_f}
%	|f(x,t)|\leq K_{f} (1+t^{2^{\star}_{Q}-1})  \quad \textrm{ for a.e. } x \in \Omega \textrm{ and for  every } t>0.
%\end{equation}
%Here and throughout, $\Omega\subseteq\mathbb{G}$ is a fixed \emph{bounded and
%	connected open set} (as in \eqref{EqProblem}).
\begin{defin}\label{def:weak_sub_super_sol}
	Let $\Omega\subseteq\mathbb{G}$ be an open, bounded and connected set.
We say that a function $u \in S^{1}_{0}(\Omega)$ is a weak subsolution (resp. supersolution) of \eqref{EqProblem} if it satisfies the following properties:
	\begin{itemize}
		\item[(i)] $u>0$ in $\Omega$.
		\item[(ii)] For every $0\leq \varphi \in C^{\infty}_{0}(\Omega)$, it holds that
		\begin{equation}\label{eq:weak_sub_super_sol}
			\int_{\Omega}\langle \nabla_{\mathbb{G}}u, \nabla_{\mathbb{G}}\varphi \rangle_{\mathfrak{g}_{1}} \leq (\textrm{resp. } \geq)  \int_{\Omega}\left(\lambda u^{q}  + \ u^{2^{\star}_{Q}}\right)\varphi.
		\end{equation}
	\end{itemize}
	Finally, we say that $u \in S^{1}_{0}(\Omega)$ is a weak solution of \eqref{EqProblem} if it is both a weak subsolution and a weak supersolution of \eqref{EqProblem} without the non-negativity condition on $\varphi$.
\end{defin}

\medskip

Let us close this section recalling a few results on the Sobolev inequality \eqref{eq:Sobolev_Ineq}.
\begin{lemma}
	The best Sobolev constant in \eqref{eq:Sobolev_Ineq} (with $\mathcal{O} = \mathbb{G}$) is achieved by a positive function $T \in S^{1}_{0}(\mathbb{G})$, and it is characterized as follows
	\begin{equation}\label{eq:def_best_constant}
		S_{\mathbb{G}} := \inf_{f\in S^{1}_{0}(\mathbb{G})} \dfrac{\| |\nabla_{\mathbb{G}}f| \|^{2}_{L^{2}(\mathbb{G})}}{	\|f\|^{2}_{L^{2_{Q}^{\star}}(\mathbb{G})}}.
	\end{equation}
	Up to a constant, $T \in S^{1}_{0}(\mathbb{G})$ is also a weak solution to 
	\begin{equation}\label{eq:Critical_PDE_in_G}
		-\Delta_{\mathbb{G}}u = u^{2^{\star}_{Q}-1} \quad \textrm{in } \mathbb{G}.
	\end{equation}
	Moreover, the following holds:
	\begin{itemize}
		\item there exists a positive constant $M_1>0$ such that 
		\begin{equation}\label{eq:Stima_Bonf_Ugu}
			T(g) \leq M_1 \, \min \{1, |g|_{\mathbb{G}}^{2-Q}\}, \quad \textrm{ for every } g \in \mathbb{G},
		\end{equation} 
		\item there exists a positive constant $M_2 >0$ such that
		\begin{equation}\label{eq:Stima_Loiudice}
			T(g) \geq M_2 \, \dfrac{|B_{1}(0)|}{(1+|g|_{\mathbb{G}})^{Q-2}}, \quad \textrm{ for every } g \in \mathbb{G}.
		\end{equation}
	\end{itemize}
	\begin{proof}
		The fact that the best Sobolev constant in \eqref{eq:Sobolev_Ineq} (with $\mathcal{O} = \mathbb{G}$) is achieved has been proved in \cite{GaroVa}. We refer to \cite[Theorem 3.4]{BoUg} for a proof of \eqref{eq:Stima_Bonf_Ugu} and to  \cite[Lemma 3.2]{Loiudice1} for a proof of \eqref{eq:Stima_Loiudice}.
	\end{proof}
\end{lemma}

As it is well known from the seminal paper \cite{BN}, a major role in finding solutions to critical problems is played by a suitable localized version of the family of minimizers of \eqref{eq:Sobolev_Ineq}. To be more precise, let $T \in S^{1}_{0}(\mathbb{G})$ be a minimizer of \eqref{eq:def_best_constant}. For every $\varepsilon >0$ define the rescaled function
\begin{equation}\label{eq:Minimi_riscalati}
	T_{\varepsilon}(g) := \varepsilon^{(2-Q)/2} T \left(\delta_{1/\varepsilon}(g)\right).
\end{equation}
Let further be $R>0$ such that $B_{R}(0)\subset \Omega$ and let $\varphi \in C^{\infty}_{0}(B_{R}(0))$ be a cut-off function such that $0\leq \varphi \leq 1$ and $\varphi \equiv 1$ in $B_{R/2}(0)$. Finally, define the family of functions 
\begin{equation}\label{eq:def_u_varepsilon}
	U_{\varepsilon}(g):= \varphi(g) T_{\varepsilon}(g), \quad g \in \mathbb{G}.
\end{equation}
Now, we have the following

\begin{lemma}\label{lem:Stime_Talentiane_modificate}
	Let $T_{\varepsilon}$ and $U_{\varepsilon}$ be as above.
	The following holds:
	\begin{itemize}
			\item[i)] Due to scaling invariance, 
		\begin{equation}
			\| |\nabla_{\mathbb{G}}T_{\varepsilon}|\|^{2}_{L^{2}(\mathbb{G})} = \| T_{\varepsilon}\|^{2^{\star}_{Q}}_{L^{2_{Q}^{\star}}(\mathbb{G})} = S_{\mathbb{G}}^{Q/2}.
		\end{equation}
		\item[ii)] The function $U_{\varepsilon}$ satisfies the following estimates as $\varepsilon \to 0^{+}$
		\begin{align}
			\| |\nabla_{\mathbb{G}}U_{\varepsilon}|\|^{2}_{L^{2}(\mathbb{G})} & =  S_{\mathbb{G}}^{Q/2} + O(\varepsilon^{Q-2}) \label{eq:Stima_grad_u_varepsilon}\\
			\|U_{\varepsilon}\|^{2^{\star}_{Q}}_{L^{2_{Q}^{\star}}(\mathbb{G})}& = S_{\mathbb{G}}^{Q/2} + O(\varepsilon^{Q}) \label{eq:Stima_u_varepsilon}. 
		\end{align}
	\end{itemize}	
	\begin{proof}
		We refer to \cite[Lemma 3.3]{Loiudice1} for a proof of both \eqref{eq:Stima_grad_u_varepsilon} and \eqref{eq:Stima_u_varepsilon}.
	\end{proof}
\end{lemma}
 
%\subsection{Strong Maximum Principle and Harnack inequality} Now we have reviewed
%the basic concepts concerning the Carnot groups setting,  
%we briefly recall here below a \emph{Weak Harnack inequality} a \emph{Strong Maximum Principle} for the operator
%$$\text{$L = -\Delta_\mathbb{G}+c(x)$ (with $c\leq 0$)}.$$
%We refer e.g. to \cite[Proposition 2.2 and Proposition 2.4]{BiGaVe} for an explicit proof.
%
%\begin{prop}[Weak Harnack inequality for $-\Delta_\mathbb{G}+c$] \label{prop:WHarnack}
% Let $\mathcal{O}\subseteq\R^n$ be a bounded open set, and let 
% $c\in L^\infty(\mathcal{O}),\,c\geq 0.$
% Moreover, let $u\in S_0^1(\mathcal{O})$ be a \emph{weak supersolution}
% of the equation 
% \begin{equation} \label{eq:PDEOrderzeroHarnack}
%  \text{$-\Delta_\mathbb{G}u+cu = 0$ in $\mathcal{O}$},
%  \end{equation}
% that is
% \emph{(}see also the subsequent Definition \ref{def:weak_sub_super_sol}\emph{)},
% $$\int_\mathcal{O}\langle\nabla_\mathbb{G}u,\nabla_\mathbb{G}\varphi\rangle_{\mathfrak{g}_{1}}
% +\int_\mathcal{O} cu\varphi\geq 0\quad\text{
% for every $\varphi\in C_0^\infty(\mathcal{O}),\,\varphi\geq 0$}.$$
% If, in addiction, $u\geq 0$ a.e.\,in $\mathcal{O}$, there exist
% constants $c_0 > 0$ and $p_0\in (0,1)$,
% \emph{independent of~$u$}, such that,
% for every $g_0\in\mathcal{O},\,r > 0$ with $B_{4r}(g_0)\Subset\mathcal{O}$, we have
% \begin{equation} \label{eq:HarnackInequ}
%  \Big(-\!\!\!\!\!\!\!\int_{B_{3r}(g_0)}|u|^{p_0}\Big)^{1/p_0}\leq c_0\,\inf_{B_{r}}u.
% \end{equation}
%\end{prop}
%As a simple consequence of Proposition \ref{prop:WHarnack}, we get the following 
%\begin{coro} \label{cor:BrezisNirenbergpernoi}
% Let $\mathcal{O}\subseteq\mathbb{G}$ be a bounded open set,
% and let $c\in L^\infty(\mathcal{O}),\,\text{$c\geq 0$}$. Furthermore, let
% $u\in S_0^1(\mathcal{O})$ be a weak superso\-lu\-tion of 
%equation \eqref{eq:PDEOrderzeroHarnack} satisfying
%$$\text{$u > 0$ a.e.\,on every open ball $B = B_r(g_0)\Subset\Omega$}.$$
%Then, for every open set $\mathcal{O}'\Subset\mathcal{O}$ there exists $C = C(\Oo',u) > 0$ such that
%$$\text{$u\geq C(\Oo',u) > 0$ a.e.\,in $\Oo'$}.$$
%\end{coro}
%
%As in the Euclidean case, Proposition \ref{prop:WHarnack} is a crucial tool to prove the following 
%\begin{prop}[Strong Maximum Principle for $L$] \label{prop:SMP}
% Let $\mathcal{O}\subseteq \mathbb{G}$ be a \emph{bounded and connected} 
% open set, and let
% $u\in S^1_0(\mathcal{O}),\,u\geq 0$ be a \emph{weak supersolution} of \eqref{eq:PDEOrderzeroHarnack}. 
% Then,
% \begin{equation} \label{eq:SMPConclusion}
%  \text{either $u \equiv 0$ or $u > 0$ a.e.\,in $\mathcal{O}$}.
% \end{equation}
%\end{prop}


%\subsection{Weak sub/supersolutions of \texorpdfstring{$\eqref{EqProblem}_{\lambda}$}{Main Problem}}
%We are now ready to properly set the definition of weak sub/supersolution of \eqref{EqProblem}$_\lambda$.
% However, we prefer to provide such definitions for the following slightly more general singular Dirichlet problem:
%\begin{equation}\label{eq:General_Singular_Problem}
%	\left\{\begin{array}{rl}
%		-\Delta_{\mathbb{G}}u = \lambda \, u^{q} + f(x,u)& \textrm{ in } \Omega,\\
%		u>0 & \textrm{ in } \Omega,\\
%		u=0 & \textrm{ on } \partial \Omega,
%	\end{array}\right.
%\end{equation}
%\noindent where $q \in (0,1)$, $\lambda >0$ and $f: \Omega \times (0,+\infty)\to \mathbb{R}$ is a Carath\'{e}odory function satisfying the following critical growth condition: there exists a positive constant $K_{f}>0$ such that
%\begin{equation}\label{eq:Growth_of_f}
%	|f(x,t)|\leq K_{f} (1+t^{2^{\star}_{Q}-1})  \quad \textrm{ for a.e. } x \in \Omega \textrm{ and for  every } t>0.
%\end{equation}
%Here and throughout, $\Omega\subseteq\mathbb{G}$ is a fixed \emph{bounded and
%connected open set} (as in \eqref{EqProblem}).
%\begin{defin}\label{def:weak_sub_super_sol}
%	Let $f: \Omega \times (0,+\infty)\to \mathbb{R}$ be a Carath\'{e}odory function satisfying the growth condition \eqref{eq:Growth_of_f}. We say that a function $u \in S^{1}_{0}(\Omega)$ is a weak subsolution (resp. supersolution) of \eqref{eq:General_Singular_Problem} if it satisfies the following properties:
%	\begin{itemize}
%		\item[(i)] $u>0$ in $\Omega$.
%		\item[(ii)] For every $0\leq \varphi \in C^{\infty}_{0}(\Omega)$, it holds that
%		\begin{equation}\label{eq:weak_sub_super_sol}
%			\int_{\Omega}\langle \nabla_{\mathbb{G}}u, \nabla_{\mathbb{G}}\varphi \rangle_{\mathfrak{g}_{1}} \leq (\textrm{resp. } \geq) \lambda \int_{\Omega}u^{q}\varphi + \int_{\Omega}f(x,u)\varphi.
%		\end{equation}
%	\end{itemize}
%	Finally, we say that $u \in S^{1}_{0}(\Omega)$ is a weak solution of \eqref{eq:General_Singular_Problem} if it is both a weak subsolution and a weak supersolution of \eqref{eq:General_Singular_Problem} without the non-negativity condition on $\varphi$.
%\end{defin}


% \begin{rmk} \label{rem:defweaksolPb}
%	We now list here below some comments concerning the above Definition \ref{def:weak_sub_super_sol}. In what follows,
%	we tacitly understand that $f:\Omega\to(0,+\infty)$ is a Carath\'eodory function satisfying
%	the growth
%	assumption \eqref{eq:Growth_of_f}.
%	\medskip
%	
%	1)\,\,If $u\in S^{1}_{0}(\Omega)$ is a weak sub/supersolution of problem
%	\eqref{eq:General_Singular_Problem}, all the integrals appearing in \eqref{eq:weak_sub_super_sol} 
%	\emph{exist and are finite}.
%	Indeed, by H\"older's inequality, for every 
%	$\varphi\in C_0^\infty(\Omega),\,\text{$\varphi\geq 0$ in $\Omega$}$, 
%	we obtain 
%	\begin{equation} \label{eq:BrhouphiFinite}
%				\left|\int_{\Omega}\langle \nabla_{\mathbb{G}}u, \nabla_{\mathbb{G}}\varphi \rangle_{\mathfrak{g}_{1}}\right|  \leq 
%			\| u\|_{S_0^1(\Omega)}\|\varphi\|_{S_0^1(\Omega)}
%	\end{equation}
%	Moreover, using \eqref{eq:Growth_of_f} and the Sobolev inequality, we also get
%	\begin{equation} \label{eq:termfFinite}
%		\begin{split}
%			\int_\Omega f(x,u)\,|\varphi|& \leq K_f\Big(\|\varphi\|_{L^1(\Omega)}
%			+ \int_{\Omega}|u|^{2^{\star}_{Q}-1}|\varphi|\,dx\Big) \\
%			& (\text{using H\"older's inequality}) \\
%			& \leq  K_f\big(\|\varphi\|_{L^1(\Omega)}
%			+\|u\|_{L^{2^{\star}_{Q}}(\Omega)}^{2^{\star}_{Q}-1}\cdot\|\varphi\|_{L^{2^{\star}_{Q}}(\Omega)}\big) \\
%			& \leq K_f\big(\|\varphi\|_{L^1(\Omega)}
%			+S_{\mathbb{G}}^{-2^{\star}_{Q}/2}\| u\|_{S_0^1(\Omega)}^{2^{\star}_{Q}-1}\cdot\| \varphi\|_{S_0^1(\Omega)}\big) \\
%			& (\text{again by H\"older's inequality and Poincaré inequality}) \\
%			& \leq C\| \varphi\|_{S_0^1(\Omega)}\big(1+\| u\|_{S_0^1(\Omega)}^{2^{\star}_{Q}-1}\big)<+\infty,
%		\end{split}
%	\end{equation}
%	where $S_{\mathbb{G}} > 0$ is the best Sobolev constant in $\mathbb{G}$, and $C > 0$ depends on $Q,\,f,\,|\Omega|$.
%	Fi\-nally, since \emph{we are assuming that $u^{-\gamma}\in L^1_{\mathrm{loc}}(\Omega)$}, we obviously
%	have
%	$$0\leq \int_\Omega u^{-\gamma}\varphi  \leq 
%	\|\varphi\|_{L^\infty(\Omega)}\int_{\mathrm{supp}(\varphi)}u^{-\gamma}< +\infty. 
%	$$
%	We explicitly stress that this last estimate 
%	(which is related with the \emph{singular term} $u^{-\gamma}$) is the unique estimate involving
%	the $L^\infty$-norm of the test function $\varphi$; on the contrary, 
%	estimates \eqref{eq:BrhouphiFinite}-\eqref{eq:termfFinite} only involve
%	the \emph{$S_0^1$-norm} of $\varphi$.
%	\vspace*{0.1cm}
%	
%	2)\,\,If $u\in S^{1}_{0}(\Omega)$ is a weak \emph{solution} of
%	problem
%	\eqref{eq:General_Singular_Problem}, and if $\varphi\in C_0^\infty(\Omega)$ is a {non-negative}
%	test function (that is, $\varphi\geq 0$ in $\Omega$), by \eqref{eq:BrhouphiFinite}-\eqref{eq:termfFinite}
%	we have
%	\begin{align*}
%		0\leq \int_\Omega u^{-\gamma}\varphi  & = 
%		\frac{1}{\lambda}\Big(\int_{\Omega}\langle \nabla_{\mathbb{G}}u, \nabla_{\mathbb{G}}\varphi \rangle_{\mathfrak{g}_{1}}
%		+ \int_\Omega f(x,u)\varphi \Big)
%		\\
%		& \leq C\| \varphi\|_{S_0^1(\Omega)}\big(1+\| u\|_{S_0^1(\Omega)}^{2^{\star}_{Q}-1}\big);
%	\end{align*}
%	from this, by using a standard \emph{density argument}, we can easily prove the following facts:
%	\begin{itemize}
%		\item[a)] $u^{-\gamma}\varphi\in L^1(\Omega)$
%		\emph{for every $\varphi\in S^{1}_{0}(\Omega)$};
%		\item[b)] identity \eqref{eq:weak_sub_super_sol}
%		actually holds \emph{for every $\varphi\in S^{1}_{0}(\Omega)$}.
%	\end{itemize}
%	
%	
%	3)\,\,Let $u\in S^{1}_{0}(\Omega)$ be a weak \emph{solution} of problem \eqref{eq:General_Singular_Problem}.
%	Since, in particular, we know that
%	$u > 0$ a.e.\,in $\Omega$ (and $u\equiv 0$ a.e.\,in $\mathbb{G}\setminus\Omega$), it is quite easy
%	to recognize that $u$ is a \emph{weak supersolution of the equation}
%	$$-\Delta_{\mathbb{G}} u  = 0\quad\text{in $\Omega$},$$
%	in the sense of Definition \ref{def:weak_sub_super_sol}.
%	As a consequence, we are entitled to apply \cite{GutLan}, ensuring that
%	\begin{equation*}
%		\begin{gathered}
%			\text{for every $\mathcal{O}\Subset\Omega$ there exists $c(\mathcal{O},u) > 0$ 
%				s.t.\,$u\geq c(\mathcal{O},u) > 0$ a.e.\,in $\mathcal{O}$}.
%		\end{gathered}
%	\end{equation*}
%\end{rmk}
%


\medskip


\subsection{Auxiliary sublinear problem}\label{subsec:Auxiliary}
A major role in what follows will be played by the weak solution of the following problem
%
%The first one concerns the 
%\emph{un\-per\-tur\-bed}, purely sublinear version of problem \eqref{EqProblem}$_\lambda$, that is,
\begin{equation}\label{eq:Sublinear_Problem}
	\left\{\begin{array}{rll}
				-\Delta_{\mathbb{G}}u &= \lambda\, u^{q} & \textrm{ in } \Omega,\\
				u&>0 & \textrm{ in } \Omega,\\
				u&=0 & \textrm{ on } \partial \Omega,
	\end{array}\right.
\end{equation}
\noindent where $q \in (0,1)$ and $\lambda >0$. We say that $u \in S^{1}_{0}(\Omega)$ is a weak solution of \eqref{eq:Sublinear_Problem} analogously to Definition \ref{def:weak_sub_super_sol}.
In this context, we have the following

\begin{teo}\label{thm:Sublinear_Problem}
	Let $\Omega\subseteq\mathbb{G}$
	be a bounded open set. Moreover, let $q \in (0,1)$ and $\lambda >0$. 
	Then, there exists a unique weak solution $\underline{u}_{\lambda} \in S^{1}_{0}(\Omega) \cap L^{\infty}(\Omega)$ to \eqref{eq:Sublinear_Problem}. Moreover, $\underline{u}_{\lambda}$ is a global minimizer in the $S^{1}_{0}(\Omega)$-topology
	of the functional 
		\begin{equation}\label{eq:def_Functional_Singular}
				J_{\lambda}(u):= \dfrac{1}{2}\int_{\Omega}|\nabla_{\mathbb{G}}u|^{2} 
				-\dfrac{\lambda}{q+1} \int_{\Omega}|u|^{q+1}.
			\end{equation}
		Finally, we also have that $J_{\lambda}(\underline u_\lambda)<0$.
   \end{teo}
The above result is actually a particular case of \cite[Theorem 1.1, Proposition 5.1]{BPV}, where Brezis-Oswald-type results have been proved for more general H\"{o}rmander operators.

\medskip


%		\begin{proof}
%%		The proof of this theorem is very similar
%%		to that of  \cite[Theorem 1.1]{SunWuLong} (for the  solvability of
%%		\eqref{eq:Singular_Problem}), and exploits a classical scheme of Stampacchia (for the
%%		global boundedness of $w_\lambda$); however, we present it
%%		here in detail for the sake of completeness.
%		
%			To ease the readability, we split the proof into five steps.
%			\medskip
%			
%			\textsc{Step I).} In this first step we prove that $J_\lambda$ possesses a 
%			\emph{global and strictly negative minimum}, which is attained by
%			a function $w_\lambda\in S_0^1(\Omega)\setminus\{0\}$ 
%			such that $w_\lambda \geq 0$ a.e.\,in $\Omega$.
%			\vspace*{0.1cm}
%
%			To this we first observe that, since $S_0^1(\Omega)$ is a Hilbert
%			space and since $0<1-\gamma<1$, the functional $J_\lambda$ is weakly
%			lower semicontinuous on $S_0^1(\Omega)$ (notice that $\||\nabla_\mathbb{G}u|\|_{L^2(\Omega)}$
%			is precisely the \emph{norm} of $S_0^1(\Omega)$). Moreover, owing to the Sobolev inequality
%			\eqref{eq:Sobolev_Ineq} and using H\"older's inequality, we have
%			\begin{equation*}
%				J_{\lambda}(u) \geq  \dfrac{\|u\|^{2}_{S^{1}_{0}(\Omega)}}{2} - 
%				\dfrac{\lambda C}{1-\gamma}\|u\|^{1-\gamma}_{S^{1}_{0}(\Omega)}
%			\end{equation*}
%			(for some positive constant $C$ only depending on $Q$ and $\Omega$),
%			and this proves that $J_\lambda$ is \emph{coercive} on $S_0^1(\Omega)$.
%			Gathering these facts, we then derive that 
%			$J_\lambda$
%			attains a \emph{global minimum}, that is, there exists 
%			a function $w_\lambda\in S_0^1(\Omega)$
%			such that
%			$$J_\lambda(w_\lambda) = \min_{u\in S_0^1(\Omega)}J_\lambda(u)=m_\lambda.$$
%			We now observe that, if $\varphi\in C_0^\infty(\Omega)$ is a fixed non-vanishing 
%			function, we have
%			$$J_\lambda(t\varphi) = \frac{t^2}{2}\int_{\Omega}|\nabla_{\mathbb{G}}\varphi|^{2} 
%			-\dfrac{\lambda t^{1-\gamma}}{1-\gamma} \int_{\Omega}|\varphi|^{1-\gamma}\to-\infty
%			\quad\text{as $t\to0^+$}$$
%			(recall that $0<1-\gamma<1$); as a consequence, we have
%			 $m_\lambda < 0$, and thus $w_\lambda\not\equiv 0$.
%
%			 Finally, taking into account that $J_\lambda(|u|)=J_\lambda(u)$ for every $u\in S_0^1(\Omega)$,
%			 by possibly replacing the function $w_\lambda$ with $|w_\lambda|$, we can assume that
%			 $w_\lambda \geq 0$ a.e.\,in $\Omega$.
%			 \medskip
%			 
%			 \textsc{Step II).} In this second step we prove that the global, non-negative
%			 minimizer $w_\lambda\in S_0^1(\Omega)$ 
%			 obtained in the previous step is actually strictly positive (a.e.) in $\Omega$.
%			 
%			 To this end, we arbitrarily fix a \emph{non-negative} function $\varphi\in C_0^\infty(\Omega)$
%			 and we notice that, since 
%			 the function $w_\lambda$ is a global minimizer of $J_\lambda$, we have
%			 \begin{align*}
%			  0 & \leq \liminf_{t\to 0^+}\Big\{
%			  \frac{1}{t}(J_\lambda(w_\lambda+t\varphi)-J_\lambda(w_\lambda))\Big\} 
%			  \leq\int_\Omega\langle\nabla_\mathbb{G}w_\lambda,
%			  \nabla_\mathbb{G}\varphi\rangle_{\mathfrak{g}_{1}},
%			 \end{align*}
%			 and thus $-\Delta_\mathbb{G}u\geq 0$ in the weak sense on $\Omega$;
%			 from this, by using the Strong Maximum Principle in Proposition 
%			 \ref{prop:SMP} (and since that $w_\lambda\not\equiv 0$), we deduce that 
%			 $w_\lambda > 0$ a.e.\,in $\Omega$.
%			 \medskip
%			 
%			 \textsc{Step III).} In this third step, we prove that the function  $w_\lambda\in S_0^1(\Omega)$
%			 (which we have already proved to be a glo\-bal and strictly positive minimzer of
%			 $J_\lambda)$ is a weak solution of 
%			 \begin{equation} \label{eq:PDEpurelysingular}
%			  -\Delta_\mathbb{G}w_\lambda = \frac{\lambda}{w_{\lambda}^\gamma}\quad\text{in $\Omega$},
%			 \end{equation}
%			 and hence of the whole problem \eqref{eq:Singular_Problem}.
%			 
%			 To this end we first observe that, since $w_\lambda$ is a global minimizer of $J_\lambda$, the point $t_0 = 0$ is a 
%			 \emph{global minimum} for the (smooth) map $h:\mathbb{R}\to\mathbb{R}$ defined by
%			 $$h(t) = J_\lambda\big((1+t)w_\lambda)\big)
%			 = \frac{(1+t)^2}{2}\int_\Omega|\nabla_\mathbb{G}w_\lambda|^2-
%			 \frac{\lambda(1+t)^{1-\gamma}}{1-\gamma}\int_\Omega|w_\lambda|^{1-\gamma};$$
%			 as a consequence, we have
%			 \begin{equation} \label{eq:operatorPartzero}
%			 	0 = h'(0) = \int_\Omega|\nabla_\mathbb{G}w_\lambda|^2-
%			 \lambda\int_\Omega|w_\lambda|^{1-\gamma}. 
%			 \end{equation}
%			 Moreover, again by the minimality of $J_\lambda(w_\lambda)$, we also get
%			 \begin{align*}
%			 	0&\leq J_\lambda(w_\lambda+t\varphi)
%			 	-J_\lambda(w_\lambda) \\
%			 	& \leq \int_\Omega\langle\nabla_\mathbb{G}w_\lambda,\nabla_\mathbb{G}\varphi\rangle_{\mathfrak{g}_{1}}
%			 + \frac{t^2}{2}\int_\Omega |\nabla_\mathbb{G}\varphi|^2
%			  - \frac{\lambda}{1-\gamma}\int_\Omega\big[
%			 (w_\lambda+t\varphi)^{1-\gamma}-w_\lambda^{1-\gamma}\big]
%			 \end{align*}
%			 for every \emph{non-negative function $\varphi\in S_0^1(\Omega)$ and every $t > 0$};
%			 thus, by letting $t\to 0^+$ with the aid of the Fatou Lemma
%			 (and by the arbitrariness of $\varphi$),
%			 we obtain 
%			 \begin{equation} \label{eq:dausareperconcludereHaitao}
%			 \int_\Omega\big(\langle\nabla_\mathbb{G}w_\lambda,\nabla_\mathbb{G}\varphi\rangle_{\mathfrak{g}_{1}}
%			 - {\lambda} w_\lambda^{-\gamma}\varphi\big)\geq 0,\quad\text{for every
%			 $\varphi\in S_0^1(\Omega),\,\varphi\geq 0$}.
%			 \end{equation}
%			 With estimates \eqref{eq:operatorPartzero}-\eqref{eq:dausareperconcludereHaitao} at hand,
%			 we can easily complete the demonstration of this step. 
%			 
%			 Indeed, given any
%			 $\psi\in S_0^1(\Omega)$, by applying
%			 \eqref{eq:dausareperconcludereHaitao} to
%			 $$\varphi = (w_\lambda+\varepsilon\psi)_+\in S_0^1(\Omega)$$
%			 (and with $\varepsilon > 0$ arbitrarily fixed), and by using \eqref{eq:operatorPartzero}, we get
%			 \begin{align*}
%			  & 0 \leq \int_\Omega\big(\langle\nabla_\mathbb{G}w_\lambda,\nabla_\mathbb{G}\varphi\rangle_{\mathfrak{g}_{1}}
%			 - {\lambda}w_\lambda^{-\gamma}\varphi\big)
%			 \\
%			 & \qquad = \int_{\{w_\lambda+\varepsilon\psi\geq 0\}}\big(\langle\nabla_\mathbb{G}w_\lambda,\nabla_\mathbb{G}(w_\lambda+\varepsilon\psi)\rangle_{\mathfrak{g}_{1}}
%			 - {\lambda}w_\lambda^{-\gamma}(w_\lambda+\varepsilon\psi)\big) \\
%			 & \qquad = \int_{\Omega}\big(\langle\nabla_\mathbb{G}w_\lambda,\nabla_\mathbb{G}(w_\lambda+\varepsilon\psi)\rangle_{\mathfrak{g}_{1}}
%			 - {\lambda}w_\lambda^{-\gamma}(w_\lambda+\varepsilon\psi)\big)	 \\
%			 & \qquad\qquad
%			 - \int_{\{w_\lambda+\varepsilon\psi< 0\}}\big(\langle\nabla_\mathbb{G}w_\lambda,\nabla_\mathbb{G}(w_\lambda+\varepsilon\psi)\rangle_{\mathfrak{g}_{1}}
%			 - {\lambda}w_\lambda^{-\gamma}(w_\lambda+\varepsilon\psi)\big)\\
%			 & \qquad
%			 = \Big(\int_\Omega|\nabla_\mathbb{G}w_\lambda|^2-
%			 \lambda\int_\Omega|w_\lambda|^{1-\gamma}\Big)
%			  +\varepsilon\int_\Omega\big(\langle\nabla_\mathbb{G}w_\lambda,\nabla_\mathbb{G}\psi\rangle_{\mathfrak{g}_{1}}
%			 - {\lambda} w_\lambda^{-\gamma}\psi\big) \\
%			 & \qquad\quad - \int_{\{w_\lambda+\varepsilon\psi< 0\}}\big(\langle\nabla_\mathbb{G}w_\lambda,\nabla_\mathbb{G}(w_\lambda+\varepsilon\psi)\rangle_{\mathfrak{g}_{1}}
%			 - {\lambda}w_\lambda^{-\gamma}(w_\lambda+\varepsilon\psi)\big)
%			 \\
%			 & \qquad (\text{here we use \eqref{eq:operatorPartzero}}) \\
%			 & \qquad = \varepsilon\int_\Omega\big(\langle\nabla_\mathbb{G}w_\lambda,\nabla_\mathbb{G}\psi\rangle_{\mathfrak{g}_{1}}
%			 - {\lambda} w_\lambda^{-\gamma}\psi\big) \\
%			 & \qquad\quad - \int_{\{w_\lambda+\varepsilon\psi< 0\}}\big(\langle\nabla_\mathbb{G}w_\lambda,\nabla_\mathbb{G}(w_\lambda+\varepsilon\psi)\rangle_{\mathfrak{g}_{1}}
%			 - {\lambda}w_\lambda^{-\gamma}(w_\lambda+\varepsilon\psi)\big)
%			 \\
%			& \qquad \leq 
%			\varepsilon\int_\Omega\big(\langle\nabla_\mathbb{G}w_\lambda,\nabla_\mathbb{G}\psi\rangle_{\mathfrak{g}_{1}}
%			 - {\lambda} w_\lambda^{-\gamma}\psi\big) 
%			 - \varepsilon\int_{\{w_\lambda+\varepsilon\psi< 0\}}\langle\nabla_\mathbb{G}w_\lambda,\nabla_\mathbb{G}\psi\rangle, 
%			  \end{align*}
%			 where we have used the fact that $w_\lambda > 0$ a.e.\,in $\Omega$. 
%			 In view of this fact, we also recognize that the Lebesgue measure of the set $A_\varepsilon = \{w_\lambda+\varepsilon\psi< 0\}$
%			 goes to $0$ as $\varepsilon\to 0^+$: indeed, since
%			 $$\psi < -\frac{1}{\varepsilon}w_\lambda < 0\quad\text{on $A_\varepsilon$},$$
%			 and since $\Omega\supseteq A_\varepsilon$ has finite measure, we have
%			 $$\lim_{\varepsilon\to 0^+}|A_\varepsilon| = 
%			 \Big|\bigcap_{\varepsilon > 0}A_\varepsilon\Big| = |\{w_\lambda\leq 0\}| = 0.$$
%			 Hence, from the above estimate
%			 (and by the arbitrariness of $\varepsilon > 0$), we obtain
%			 \begin{align*}
%			 	\int_\Omega\big(\langle\nabla_\mathbb{G}w_\lambda,\nabla_\mathbb{G}\psi\rangle_{\mathfrak{g}_{1}}
%			 - {\lambda} w_\lambda^{-\gamma}\psi\big) 
%			\geq \lim_{\varepsilon\to 0^+}\int_{\{w_\lambda+\varepsilon\psi< 0\}}\langle\nabla_\mathbb{G}w_\lambda,\nabla_\mathbb{G}\psi\rangle = 0.
%			 \end{align*}
%			 This, together with the arbitrariness of
%			 $\psi \in S_0^1(\Omega)$, allows us to finally conclude that 
%			 $w_\lambda$ is a weak solution of equation
%			 \eqref{eq:PDEpurelysingular}, as desired.
%			 			 \medskip
%			 
%			 \textsc{Step IV).} In this fourth step we show that $w_\lambda\in L^\infty(\Omega)$.
%			 To this end we first observe that, since we \emph{already know} that
%			 $w_\lambda$ is a weak solution
%			 of \eqref{eq:Singular_Problem}, for every $k> 1$ we have
%			 $$
%			 \int_{\Omega}
%			  \langle\nabla_\mathbb{G}w_\lambda,\nabla_\mathbb{G}[(w_\lambda-k)_+]\rangle_{\mathfrak{g}_{1}}
%			  -\lambda\int_{\Omega}w_\lambda^{-\gamma}(w_\lambda-k)_+ = 0;$$
%			  from this, by using the H\"older inequality and the weighted Young inequality, together
%			  with the continuous embedding of $S_0^1(\Omega)$ into $L^{2_Q^\star}(\Omega)$, we get
%	\begin{align*}
%	 \int_{\Omega}
%			  |\nabla_\mathbb{G}[(w_\lambda-k)_+]|^2 & \leq
%			   \frac{\lambda}{k^\gamma}\int_{A_k}(w_\lambda-k)_+
%		\leq \frac{\lambda}{k^\gamma}|A_k|^{1-1/2_Q^\star}
%		\Big(\int_{\Omega}|(w_\lambda-k)_+|^{2_Q^\star}\Big)^{1/2_Q^\star} \\
%		& \leq c|A_k|^{1-1/2_Q^\star}\Big(\int_{\Omega}|\nabla_\mathbb{G}[w_\lambda-k)_+]|^{2}\Big)^{1/2}
%		\\
%		& \leq
%		c(\varepsilon) |A_k|^{2(1-1/2_Q^\star)}+\varepsilon\int_{\Omega}|\nabla_\mathbb{G}[w_\lambda-k)_+]|^{2},
%	\end{align*}
%	where $c(\varepsilon) > 0$ is a suitable constant depending on $\varepsilon$, and
%	$$A_k = \{w_\lambda\geq k\}.$$
%	From this, by choosing $\varepsilon > 0$ sufficiently small, we get
%	\begin{equation} \label{eq:estimLinfDaUsare}
%	 \int_{\Omega}
%			  |\nabla_\mathbb{G}[(w_\lambda-k)_+]|^2
%			  \leq c |A_k|^{2(1-1/2_Q^\star)}.
%	\end{equation}
%	With estimate \eqref{eq:estimLinfDaUsare} at hand, 
%	we can easily complete the proof of this step (and of the  whole
%	theorem). Indeed, owing to the cited \eqref{eq:estimLinfDaUsare}, for every
%	$1< k < h$ we have
%	\begin{align*}
%	 (h-k)^{2}|A_h|^{2/2_Q^\star}
%	 & \Big(\int_{A_h}(w_\lambda-k)_+^{2_Q^\star}\Big)^{2/2_Q^\star}
%	 \leq \Big(\int_{A_k}(w_\lambda-k)_+^{2_Q^\star}\Big)^{2/2_Q^\star} \\
%	 & (\text{here we use once again \eqref{eq:Sobolev_Ineq}}) \\
%	 & \leq c\int_{\Omega}
%			  |\nabla_\mathbb{G}[(w_\lambda-k)_+]|^2\leq 
%			  c |A_k|^{2(1-1/2_Q^\star)},
%	\end{align*}
%	for some constant $c > 0$ independent of $k,h$; as a consequence, we get
%	$$|A_h|\leq \frac{c}{(h-k)^{2_Q^\star}}|A_k|^{2_Q^\star-1}\quad\text{for all $1< k < h$}.$$
%	Since, obviously, $q = 2_Q^\star-1 > 1$, we can apply \cite[Lemma B.1]{KS}: this gives
%	$$|A_{d}| = |\{w_\lambda\geq d\}| = 0$$
%	for some $d > 0$, and thus $w_\lambda\in L^\infty(\Omega)$, as desired.
%	\medskip
%	
%	\textsc{Step V)}. In this last step, we prove that $w_\lambda$ is
%	the \emph{unique} weak solution of problem \eqref{eq:Singular_Problem}. To this end, let us suppose that $z_\lambda\in S_0^1(\Omega)$ is another
%	weak solution of the same problem; we then choose a smooth function $\theta\in C^\infty(\mathbb{R})$ satisfying the following properties
%		\begin{itemize}
%		 \item  $\theta(t) = 0$ for $t \leq 0$ and $\theta(t) =1$ for $t \geq 1$;
%		 \item $\theta$ is non-decreasing on $\mathbb{R}$;
%		 \end{itemize}
%		and we define (for every $\varepsilon > 0$)
%$$\theta_\varepsilon(t): = \theta\left(\frac{t}{\varepsilon}\right).$$ 
%Now, since $\varphi_\varepsilon\in S_0^1(\Omega)$, and since
%$w_\lambda,z_\lambda$ solve \eqref{eq:Singular_Problem}, we get	
%\begin{align*}
%	(\ast)\,\,&\int_{\Omega}\langle\nabla_{\mathbb{G}} w_{\lambda}, \nabla_{\mathbb{G}}(w_{\lambda}-z_\lambda)\rangle_{\mathfrak{g}_{1}} \, \theta'_{\varepsilon}(w_{\lambda}-z_\lambda)
%			- \lambda \int_{\Omega}\dfrac{\theta_{\varepsilon}(w_{\lambda}-z_\lambda)}{w_{\lambda}^{\gamma}} = 0; \\[0.1cm]
%	(\ast)\,\,& \int_{\Omega}\langle\nabla_{\mathbb{G}} z_\lambda, \nabla_{\mathbb{G}}(w_{\lambda}-z_\lambda)\rangle_{\mathfrak{g}_{1}}\, \theta'_{\varepsilon}(w_{\lambda}-z_\lambda) 
%	- \lambda \int_{\Omega}\dfrac{\theta_{\varepsilon}(w_{\lambda}-z_\lambda)}{z_\lambda^{\gamma}} = 0.	
%		\end{align*}
%	As a consequence, by subtracting the above identities, we obtain
%	\begin{equation}\label{eq:SubtractionUniqueness}
%		\begin{aligned}
%			0 &\geq  -\int_{\Omega}|\nabla(z_{\lambda}-w_\lambda)|_{\mathbb{G}}^2 \, \theta'_{\varepsilon}(w_{\lambda}-z_\lambda)
%			 \\
%			 & =  \lambda\int_\Omega\Big(\frac{1}{z_\lambda^\gamma}-\frac{1}{w_\lambda^\gamma}\Big)\theta_\varepsilon(w_\lambda-z_\lambda).
%		\end{aligned}
%	\end{equation}
%	From this, by letting $\varepsilon \to 0^+$ we derive that	
%	$$\lim_{\varepsilon\to 0^+}\int_{\Omega}\left( \dfrac{1}{z_\lambda^{\gamma}} - \dfrac{1}{w_{\lambda}^{\gamma}}\right) \theta_{\varepsilon}(w_{\lambda}-z_\lambda)
%	= \int_{\{w_{\lambda}> z_\lambda\}}\left( \dfrac{1}{z_\lambda^{\gamma}}- \dfrac{1}{w_{\lambda}^{\gamma}}\right)\leq 0,$$
%	 from which we derive that
%	$|\{ w_{\lambda}> z_\lambda\}| = 0$, that is,
%	$w_\lambda\leq z_\lambda$ a.e.\,on $\Omega$.	
%	Finally, by interchan\-ging the role of $z_\lambda$ and $w_\lambda$,
%	we conclude that
%	$$\text{$w_\lambda\equiv z_\lambda$ a.e.\,in $\Omega$},$$
%	and this proves the uniqueness of $w_\lambda$.
%			\end{proof}


We now state an {\em ad hoc} comparison principle for weak super and subsolutions of the model purely sublinear problem \eqref{eq:Sublinear_Problem}. This can be seen as a particulr instance, in the Carnot group setting, of \cite[Lemma 3.3]{ABC}, which in turn was inspired by \cite{BrezisKamin}. We refer to \cite{Ruzhansky_Suragan} for more general result whose proofs rely on the validity of suitable Picone-type inequalities.

	\begin{lemma}\label{lem:Weak_Comparison_Model}
		%Let $f$ be a Charathéodory function such that $\tfrac{f(s)}{s}$ is strictly decreasing for $s>0$. 
		Let $\lambda >0$, $q\in (0,1)$ and $v,w\in S^1_0(\Omega)$ weakly satisfy
		\begin{equation}\label{eq:per_v}
		   \left\{\begin{array}{rll}
		   	-\Delta_\G v &\leq \lambda  v^q & \textrm{ in }\Omega,\\
		    v& >0 & \textrm{ in } \Omega,\\
		    v&=0 & \textrm{ on } \partial \Omega,
		   \end{array}
		   \right.
		\end{equation}
		\noindent and 
			\begin{equation}\label{eq:per_w}
			\left\{\begin{array}{rll}
				-\Delta_\G w& \geq \lambda w^q & \textrm{in }\Omega,\\
				w &>0 & \textrm{in }\Omega,\\
				w&=0 & \textrm{on } \partial \Omega.
			\end{array}
			\right.
		\end{equation}
		 Then $w\geq v$ in $\Omega$.
		\begin{proof}
			We closely follow \cite[Proof of Lemma 3.3]{ABC}. 
			We choose first a smooth function $\theta\in C^\infty(\mathbb{R})$ satisfying the following properties
			\begin{itemize}
				\item  $\theta(t) = 0$ for $t \leq 0$ and $\theta(t) =1$ for $t \geq 1$;
				\item $\theta$ is non-decreasing on $\mathbb{R}$;
			\end{itemize}
			and we define (for every $\varepsilon > 0$)
			$$\theta_\varepsilon(t): = \theta\left(\frac{t}{\varepsilon}\right) \in S_0^1(\Omega).$$ 
			We further consider the variational formulation of both \eqref{eq:per_v} and \eqref{eq:per_w}, namely 
			\begin{equation}\label{eq:variational_for_w}
				\int_{\Omega}\langle \nabla_{\mathbb{G}}w, \nabla_{\mathbb{G}}\varphi \rangle_{\mathfrak{g}_1} \geq \lambda \int_{\Omega} w^q\varphi \quad 0\leq \varphi \in S^{1}_{0}(\Omega),
			\end{equation}
			\noindent and 
			\begin{equation}\label{eq:variational_for_v}
				\int_{\Omega}\langle \nabla_{\mathbb{G}}v, \nabla_{\mathbb{G}}\varphi \rangle_{\mathfrak{g}_1} \geq \lambda \int_{\Omega}v^q \varphi \quad 0\leq \varphi \in S^{1}_{0}(\Omega).
			\end{equation}
			We then test \eqref{eq:variational_for_w} with $\theta_{\varepsilon}(v-w) v$ and \eqref{eq:variational_for_v} with $\theta_{\varepsilon}(v-w) w$. Finally, we subtract the latter to the former, getting
			\begin{equation}\label{eq:Testata_con_theta}
			\begin{aligned}
			\int_{\Omega}&\left(w^{q-1}-v^{q-1} \right)\theta_{\varepsilon}(v-w)vw \\
			&= \int_{\Omega}\langle \nabla_{\mathbb{G}}w, \nabla_{\mathbb{G}}(v-w)\rangle_{\mathfrak{g}_1}v \, \theta'_{\varepsilon}(v-w) - \int_{\Omega}\langle \nabla_{\mathbb{G}}v, \nabla_{\mathbb{G}}(v-w)\rangle_{\mathfrak{g}_1}w \, \theta'_{\varepsilon}(v-w)\\
			&\leq \int_{\Omega}\langle \nabla_{\mathbb{G}}v,\nabla_{\mathbb{G}}(v-w)\rangle_{\mathfrak{g}_1} (v-w)\theta'_{\varepsilon}(v-w) \\
			&= \int_{\Omega}\langle \nabla_{\mathbb{G}}v, \nabla_{\mathbb{G}}(\gamma_{\varepsilon}(v-w))\rangle_{\mathfrak{g}_1} = \int_{\Omega}(-\Delta_{\mathbb{G}}v) \, \gamma_{\varepsilon}(v-w),
		\end{aligned}
		\end{equation}
			\noindent where 
			\begin{equation*}
				\gamma_{\varepsilon}(t):= \int_{0}^{t}s\theta'_{\varepsilon}(s)\, ds.
			\end{equation*}
			By construction of $\theta_{\varepsilon}$, it follows that $0\leq \gamma_{\varepsilon}(t)\leq \varepsilon$ for every $t\in \mathbb{R}$. Therefore, exploiting both \eqref{eq:Testata_con_theta}, \eqref{eq:per_v} and H\"{o}lder inequality, we find that
			\begin{equation}
					\int_{\Omega}\left(w^{q-1}-v^{q-1} \right)\theta_{\varepsilon}(v-w)vw \leq \lambda \, \varepsilon \int_{\Omega}v^q \leq C(\lambda, \|v\|_{L^{1}(\Omega)},|\Omega|,q)\,  \varepsilon.
			\end{equation}
			By letting $\varepsilon \to 0^+$, we find that
			\begin{equation*}
				\int_{\{v>w\}}\left(w^{q-1}-v^{q-1} \right)vw \leq 0,
			\end{equation*}
			\noindent from which, recalling that $q \in (0,1)$, we conclude that $|\{v>w\}|=0$. This closes the proof.
			\end{proof}
	\end{lemma}		
			
			
			
%			
%			\color{orange}
%			\begin{prop}[Proposition 5 ADV]\label{prop5adv}
%				If there exists $x_0\in \Omega$ st $v(x_0)=u(x_0)$ then $v\equiv u$ on the whole $\Omega$.
%			\end{prop}
%			\begin{teo}[questo si può fare adattando Brezis-Oswald]\label{brezisoswald_exunique}
%				If $f(x,u)$ is a Charathéodory function such that...
%				
%				Then, there exists unique the solution of the following
%				\[
%				\left\{
%				\begin{array}{rll}
%					-\Delta_\G u&= f(x,u) &\m{in }\Omega\\
%					u &\geq0 &\m{in }\Omega\\
%					u&=0 &\m{on }\de\Omega.
%				\end{array}
%				\right.
%				\]
%			\end{teo}
%			
%%			\begin{lemma}[Lemma 2.8 ADV]\label{lemma2.8adv}
%%				Consider $(v_n)_n$ a squence of non-negative functions, bounded in $S^1_0(\Omega)$ and
%%				converging weakly
%%				\[
%%				v_n\rightharpoonup v\ \m{in }S^1_0(\Omega)\quad \m{with }v_n\leq v.
%%				\]
%%				If $-\Delta_\G v_n\geq0$, then $v_n\to v$ converges strongly in $S^1_0(\Omega)$.\newline
%%			\end{lemma}
%			\color{black}

We now state an extension to the Carnot group setting of \cite[Theorem 2.4]{Struwe}. We omit the proof recalling that it is a simplified version of \cite[Lemma 3.4]{BiGaVe} where the considered functional was not differentiable i n $S_{0}^{1}(\Omega)$.

\begin{lemma}\label{lem:Perron}
	Let $\underline{u}, \overline{u}\in S^{1}_{0}(\Omega)$ be a weak subsolution and a weak supersolution, respectively, of problem \eqref{EqProblem}. 
	We assume that 
	\begin{itemize}
		\item[a)] $\underline{u}(g) \leq \overline{u}(g)$ for a.e.\,$g\in \Omega$;
		\item[b)] for every open set $\Oo\Subset\Omega$ there exists
		$C = C(\Oo,\underline{u}) > 0$ such that
		$$\text{$\underline{u}\geq C$ a.e.\,in $\Oo$}.$$
	\end{itemize}
	Then, there exists a weak solution $u \in S^{1}_{0}(\Omega)$ of \eqref{EqProblem} such that 
	$$\text{$\underline{u}(g) \leq u(g) \leq \overline{u}(g)$ for a.e. $g \in \Omega$}.$$
\end{lemma}



\section{Proof of Theorem \ref{thm:main} - Part A) and part B)}\label{sec:First_Solution}
%
%\color{red}
%\[
%I_\lambda(u)=\frac{1}{2}\int_\Omega\| |\nabla_Gu|\|^2-\frac{\lambda}{q+1}\int_\Omega u^{q+1}-\frac{1}{p+1}\int_\Omega u^{p+1}
%\]
%\color{black}
The goal of this Section is to prove the existence of a positive weak solution to \eqref{EqProblem}. 
To begin with, we define
\begin{equation}\label{eq:DefinitionLambda}
	\Lambda := \sup \{ \lambda >0: \eqref{EqProblem} \textrm{ admits a weak solution}\}.
\end{equation}
Our task now is pretty standard and it consists in the following steps: 
\begin{itemize}
	\item[I)] prove that $0<\Lambda<+\infty$;
	\item[II)] prove that problem \eqref{EqProblem} admits a weak solution for every $0<\lambda\leq \Lambda$.
\end{itemize}

\noindent We split the proof of I) in two lemmas. 
\begin{lemma}\label{lem:lambda0}
	Let $\Lambda$ as defined in \eqref{eq:DefinitionLambda}. Then $\Lambda >0$.
	\begin{proof}
		We will show that there exists a sufficiently small $\lambda >0$ such that 
		\eqref{EqProblem} has a solution. To this aim, we will use Lemma \ref{lem:Perron} exhibiting both a super and a subsolution.
		Looking for a supersolution, we consider the following auxiliary torsion problem
		\begin{equation}
			\left\{\begin{array}{rl}
				-\Delta_{\mathbb{G}} V = 1 & \textrm{in } \Omega,\\
				V = 0 & \textrm{on } \partial \Omega,
			\end{array}\right.
		\end{equation}
		\noindent whose unique solution is provided by Lax-Milgram Theorem. Moreover, by a classical Stampacchia iteration method, it holds that $V\in L^{\infty}(\Omega)$. Observe further that for any positive constant $C$,
		the function $C\cdot x^p-x$, with $p>1$, has negative value for $x>0$ sufficiently small.
		Therefore, for every $C,C'\in \R^+$
		there exists $\lambda^*>0$ such that 
		for every $\lambda<\lambda^*$, 
		\[
		\exists \, m_\lambda\in \R^+:\quad \lambda\cdot  C'\cdot m_\lambda^q+C\cdot m_\lambda^p-m_\lambda\leq0.
		\]
		We fix $\lambda<\lambda^*$ and set $C=\|V\|_{L^{\infty}(\Omega)}^p$, $C'=\|V\|_{L^{\infty}(\Omega)}^q$. 
		We define $\overline u_1:=m_\lambda V$, which weakly verifies
		\[
		\left\{\begin{array}{rll}
			-\Delta_{\G}\overline u_1&=m_\lambda \geq \lambda\overline u_1^q+\overline u_1^p &\ \m{in }\Omega\\
			\overline u_1 &>0 &\ \m{in }\Omega\\
			\overline u_1 &=0 &\ \m{on }\de\Omega,
			\end{array}\right.
		\]
				therefore it is a weak supersolution of \eqref{EqProblem}.\\
				Regarding the weak subsolution to \eqref{EqProblem}, we choose the unique solution $\underline u_\lambda$ to \eqref{eq:Sublinear_Problem}. 
				We can now conclude the proof by appealing Lemma \ref{lem:Perron}. Indeed, by Lemma \ref{lem:Weak_Comparison_Model} with $w=\overline{u}_1$ and $v=\underline{u}_{\lambda}$, we get that 
				$\overline{u}_1 \geq \underline{u}_{\lambda}$, which is condition a) of Lemma \ref{lem:Perron}. Regarding b) of Lemma \ref{lem:Perron} it is enough to recall \cite[Corollary 2.3]{BiGaVe}. This closes the proof.				
	\end{proof}
\end{lemma}

\begin{lemma}\label{lem:LambdaInf}
	Let $\Lambda$ as defined in \eqref{eq:DefinitionLambda}. Then, $\Lambda < +\infty$.
\end{lemma}
	\begin{proof}
	We consider the first eigenfunction $e_1$ of the operator $-\Delta_\G$
	with respect to the first Dirichlet eigenvalue $\mu_1$. In particular, the following characterization holds:
	\[
	\mu_1=\min\left\{\| |\nabla_\G u|\|^2_{L^2(\Omega)}:\ u\in S^1_0(\Omega)\m{ and }\|u\|^2_{L^2(\Omega)}=1\right\}.
	\]
	Knowing that $\|e_1\|_{L^2}=1$, $e_1>0$ a.e.~in $\Omega$ and $\| |\nabla_\G e_1|\|^2_{L^2}=\mu_1$,
	we have that any solution $u$ to \eqref{EqProblem} for some $\lambda$
	must verify
	\[
	\int_\Omega \langle \nabla_\G u,\nabla_\G e_1\rangle =\mu_1 \int_\Omega u e_1=\int_\Omega \lambda u^qe_1+u^pe_1.
	\]
	As $q\in (0,1)$ and $p>1$, for $\Lambda^*$ sufficiently big, we have
	\[
	\Lambda^*x^q+x^p>\mu_1x\quad\forall x\in \R^+.
	\]
	Therefore we must have $\lambda <\Lambda^*$ and this  proves $\Lambda\leq \Lambda^*<+\infty$.
\end{proof}

\medskip

%This proves part a), we prove now part b),
%namely for every $\lambda\in (0,\Lambda)$ there exists
%(at least) one weak solution of problem \eqref{EqProblem}$_\lambda$.
%%To this end, following \cite{Haitao}
%%we first establish some  preliminary results.

%\begin{lemma}\label{existencefirstsol}
%	If for $\lambda_1>0$ the problem \eqref{EqProblem}$_{\lambda_1}$ has a weak solution,
%	then for every $0<\lambda<\lambda_1$ the problem \eqref{EqProblem}$_{\lambda}$ has
%	(at least) one weak solution too.
%\end{lemma}
%\begin{proof}
%	We denote by $v_{\lambda_1}$ the solution of \eqref{EqProblem}$_{\lambda_1}$ that exists
%	by hypothesis. Therefore $v_{\lambda_1}$ is a supersolution of \eqref{EqProblem}$_{\lambda}$ for
%	any $0<\lambda<\lambda_1$. We also consider $ \underline u^*$,
%	 the unique solution to \eqref{eq:Sublinear_Problem}. Then $\underline u^*$ is a subsolution to \eqref{EqProblem}$_{\lambda_2}$.
%	
%	We verify that we can apply the Perron's methode detailed in \cite[Theorem 2.4]{Struwe}.
%	Indeed, $v_{\lambda_1}\in L^\infty(\Omega)$ and by the comparison Lemma \ref{lem:Weak_Comparison_Model}
%	we also have $v_{\lambda_1}\geq \underline u^*$. Moreover, $\underline u^*\geq0$.
%	Then, a (weak) solution $v_\lambda$ o \eqref{EqProblem}$_\lambda$ exists such that
%	\[
%	\underline u^*\leq v_\lambda\leq v_{\lambda_1}.
%	\]
%	\end{proof}
%	\color{black}
%
%\color{blue}
%In the following we prove the existence of a minimal solution $u_\lambda$
%among
%the solutions of \eqref{EqProblem}$_\lambda$, and
%we prove that $I_\lambda(u_\lambda)$ is negative.
%
%\begin{lemma}\label{lem:minsol}
%	For every $\lambda\in (0,\Lambda)$ the problem \eqref{EqProblem}$_\lambda$
%	has a minimal solution $u_\lambda$ such that $I_\lambda(u_\lambda)<0$ and for $0<\lambda_2<\lambda_1<\Lambda$
%	we have
%	$u_{\lambda_2}<u_{\lambda_1}$ on the whole $\Omega$.
%\end{lemma}
%\begin{proof}
%	We build the  sequence $(u_n)_n$ by recurrence,
%	starting from $u_0:=\underline u^*$, that is the unique 
%	solution of \eqref{eq:subsol}, and posing
%	\[
%	\left\{
%	\begin{array}{rll}
%		-\Delta_\G u_n&=\lambda u_{n-1}^q+u_{n-1}^p &\m{in }\Omega\\
%		u_n&\geq0 &\m{in }\Omega\\
%		u_n&=0 &\m{on }\de\Omega.
%		\end{array}
%	\right.
%	\]
%%	Observe that the $u_n$ is uniquely determined by the system
%%	above as a direct consequence of Theorem \ref{brezisoswald_exunique}.
% 	Using again  \cite[Theorem 1.1]{BPV} we have that $u_n$ exists unique for every $n\geq0$.
% 	We prove by recurrence that, if $v_\lambda$ is a solution of \eqref{EqProblem}$_\lambda$
% 	(such a solution exists because of Lemma \ref{existencefirstsol}), then $v_\lambda\geq u_n\geq u_{n-1}$ for every $n$.
% 	
% 	First observe that $-\Delta_\G u_1=\lambda u_0^q+u_0^p\geq-\Delta_\G u_0$, therefore
% 	by maximum principle $u_1\geq u_0$. In general for any $n\geq1$,
% 	\[
% 	-\Delta_\G u_n=\lambda u_{n-1}^q+u_{n-1}^p\geq u_{n-2}^q+u_{n-2}^p=-\Delta_\G u_{n-1},
% 	\]
% 	and so $u_{n}\geq u_{n-1}.$
% 	
% 	Now observe that $u_0\leq v_\lambda$ because of the comparison Lemma \ref{lem:Weak_Comparison_Model},
% 	after observing that $-\Delta_\G v_\lambda=\lambda v_\lambda^q+v_\lambda^p\geq\lambda v_\lambda^q$.
% 	For any $n\geq1$, 
% 	\[
% 	-\Delta_\G u_n=\lambda u_{n-1}^q+u_{n-1}^p\leq \lambda v_\lambda^q+v_\lambda^p=-\Delta_\G v_\lambda,
% 	\]
% 	therefore we obtain $u_n\leq v_\lambda$ again by maximum principle POSSO USARLO O DEVO AVERE DELLE ACCORTEZZE?
% 	
%	Therefore by Proposition \ref{prop5adv}, for every $n\geq0$
%	and for every solution $v_\lambda$ of \eqref{EqProblem}$_\lambda$,
%	\[
%	\underline u^*< u_n<v_\lambda.
%	\]
%	As a direct consequence, in this setting $\|u_n\|$ is bounded by $\| v_\lambda\|$,
%	thus there exists $u_\lambda\in S^1_0(\Omega)$ such that 
%	$u_n\rightharpoonup u_\lambda$ up to a sub-sequence.
%	Moreover, NON SONO SICURO DEL PERCHE
%	from $-\Delta_\G u_n\geq0$ we get that $u_n\to u_\lambda$
%	strongly ad $u_\lambda\leq v_\lambda$.
%	As $v_\lambda$ was chose arbitrarily among the (weak)
%	solutions of \eqref{EqProblem}$_\lambda$, this
%	proves that $u_\lambda$ is a minimal solution.
%	
%	Finally, reasoning as above we get $u_{\lambda_1}<u_{\lambda_2}$ in $\Omega$
%	if $\lambda_1<\lambda_2$. Indeed, we build the sequence $(u_n)$ as above
%	for $\lambda=\lambda_1$, then we get $u_{\lambda_2}\geq u_0$ because of the comparison
%	Lemma \ref{lem:Weak_Comparison_Model}. Inductively, this leads to $u_{\lambda_2}\geq u_n$
%	and ultimately to $u_{\lambda_2}\geq u_{\lambda_1}$.
%	 Proposition \ref{prop5adv} gives the strict inequality.
%\end{proof}
%	
%\begin{lemma}
%	When $\lambda=\Lambda$ the problem \eqref{EqProblem}$_\Lambda$ has (at least) one solution $u_\Lambda$.
%\end{lemma}	
%\begin{proof}
%	Consider an increasing sequence $(\lambda_n)\subset (0,\Lambda)$ such that $\lambda_n\to \Lambda$ for $n\to+\infty$.
%	For every $\lambda_n$, we denote by $u_n:=u_{\lambda_n}$ the minimal solution
%	of \eqref{EqProblem}$_{\lambda_n}$, that exists as proved in Lemma \ref{lem:minsol}.
%	
%	We proved that $u_n$ is an increasing sequence and $I_{\lambda_n}(u_n)<0$ for every $n$.
%	Moreover, $I_{\lambda_n}$ is Frechet differentiable at $u_n$ and $DI_{\lambda_n}(u_n)=0$.
%	Therefore, we can evaluate the differential at $u_n$ obtaining POSSIAMO DAVVERO?
%	\begin{align*}
%		0&> I_{\lambda_n}(u_n)-\frac{1}{p+1}DI_{\lambda_n}(u_n)(u_n)\\
%		&\m{SECONDO ME QUI C'è $=$ ma ADV METTON $\geq$}\\
%		&=\left(\frac{1}{2}-\frac{1}{p+1}\right)\|u_n\|^2_{S^1_0(\Omega)}+\lambda_n\left(\frac{1}{p+1}-\frac{1}{q+1}\right)\|u_n\|^{q+1}_{q+1}\\
%		&\geq\left(\frac{1}{2}-\frac{1}{p+1}\right)\|u_n\|^2_{S^1_0(\Omega)}-\lambda_n\left(\frac{1}{q+1}-\frac{1}{p+1}\right)\|u_n\|^{q+1}_{S^1_0(\Omega)}.
%	\end{align*}
%	Then (BY HOLDER?) $\|u_n\|$ is bounded in $S^1_0(\Omega)$ 
%	and $u_n\rightharpoonup u^*$ for some $u^*\in S^1_0(\Omega)$.
%	Then by $-\Delta_\G u_n\geq0$ and Lemma \ref{lemma2.8adv}
%	we get $u_n\to u^*$ strongly in $S^1_0(\Omega)$. Hence
%	$u^*$ is a solution of~\eqref{EqProblem}$_\Lambda$ (NON HO CAPITO CONCLUSIONE).
%\end{proof}

Combining Lemma \ref{lem:lambda0} with Lemma \ref{lem:LambdaInf} we get I). Let us now turn our attention to the proof of II). Firstly, let us define the functional $I_{\lambda}$ naturally associated to \eqref{EqProblem}:
\begin{equation}\label{eq:Def_Ilambda}
	I_{\lambda}(u) := \dfrac{1}{2}\int_{\Omega}|\nabla_{\G}u|^2 - \dfrac{\lambda}{q+1}\int_{\Omega}|u|^{q+1} - \dfrac{1}{2^{\star}_{Q}}\int_{\Omega}|u|^{2^{\star_{Q}}}, \quad u\in S^{1}_{0}(\Omega).
\end{equation}

\begin{teo}\label{thm:Existence_First}
	Problem \eqref{EqProblem} admits at least
	one weak solution $u_\lambda\in S^{1}_{0}(\Omega)$ for every $\lambda \in (0,\Lambda]$.
\end{teo}
\begin{proof}
	As for the proof of Lemma \ref{lem:lambda0}, we find either a weak subsolution and a weak supersolution and then apply Lemma \ref{lem:Perron}. \\
	As long as the weak subsolution is concerned, we can take the unique solution $\underline{u}_{\lambda}$ of \eqref{eq:Sublinear_Problem}. 
%	By Theorem \ref{thm:Sublinear_Problem}, we know that for every $\lambda \in (0,\Lambda)$ (and actually for all $\lambda >0$) there exists a unique solution $w_{\lambda}$ of \eqref{eq:Sublinear_Problem}, which is the Euler-Lagrange equation naturally associated with the functional $J_{\lambda}$ defined in \eqref{eq:def_Functional_Singular}.	
%	The function $w_{\lambda}$ is a weak subsolution of \eqref{eq:Main_Problem}$_\lambda$. 
	
	Let us now look for a weak supersolution. In doing this we profit of the very definition of $\Lambda$, which guarantess the existence of $\lambda' \in (\lambda, \Lambda)$ such that \eqref{EqProblem} (witht $\lambda =\lambda'$) admits a weak solution $u_{\lambda'}$. Clearly, this is a weak supersolution of \eqref{EqProblem}.\\
	By Lemma \ref{lem:Weak_Comparison_Model}, with $w = u_{\lambda'}$ and $v=w_{\lambda}$, it follows that
	\begin{equation}\label{eq:Claim}
		w_{\lambda}(g) \leq u_{\lambda'}(g), \quad \textrm{for a.e. } g\in \Omega.
	\end{equation}
%	To this aim, we proceed essentially as in the proof
%	of Theorem \ref{thm:Singular_Problem}, \textsc{Step V).} First of all, let us consider a smooth non-decreasing function $\theta: \mathbb{R}\to \mathbb{R}$ such that
%	$$\theta(t) =1 \, \textrm{ for } t \geq 1 \quad \textrm{ and } \quad \theta(t) = 0 \, \textrm{ for } t \leq 0,$$
%	\noindent and it is linked in a smooth way for $t \in (0,1)$. We further define the function
%	$$\theta_{\varepsilon}(t):= \theta \left(\dfrac{t}{\varepsilon}\right), \quad \varepsilon>0, \, t\in \mathbb{R}.$$
%	Due to its definition, we are entitled to use the function $\theta_{\varepsilon}(w_{\lambda}- u_{\lambda'})$ as a test function in both \eqref{eq:Main_Problem}$_{\lambda'}$ (solved by $u_{\lambda'}$) and \eqref{eq:Singular_Problem} (solved by $w_{\lambda}$). Thus, we have
%	\begin{equation}\label{eq:solvedBywlambda}
%		\int_{\Omega}\langle\nabla_{\mathbb{G}} w_{\lambda}, \nabla_{\mathbb{G}}(w_{\lambda}-u_{\lambda'})\rangle_{\mathfrak{g}_{1}} \, \theta'_{\varepsilon}(w_{\lambda}-u_{\lambda'})
%		- \lambda \int_{\Omega}\dfrac{\theta_{\varepsilon}(w_{\lambda}-u_{\lambda'})}{w_{\lambda}^{\gamma}} = 0,
%	\end{equation}
%	\noindent and
%	\begin{equation}\label{eq:solvedByuOverlinelambda}
%		\begin{aligned}
%			&\int_{\Omega}\langle\nabla_{\mathbb{G}} u_{\lambda'}, \nabla_{\mathbb{G}}(w_{\lambda}-u_{\lambda'})\rangle_{\mathfrak{g}_{1}}\, \theta'_{\varepsilon}(w_{\lambda}-u_{\lambda'})\\
%			& \qquad - \lambda' \int_{\Omega}\dfrac{\theta_{\varepsilon}(w_{\lambda}-u_{\lambda'})}{u_{\lambda'}^{\gamma}} - \int_{\Omega}u_{\lambda'}^{2^{\star}_{Q}-1}\theta_{\varepsilon}(w_{\lambda}-u_{\lambda'}) = 0.
%		\end{aligned}
%	\end{equation}
%	Subtracting \eqref{eq:solvedBywlambda} from  \eqref{eq:solvedByuOverlinelambda} we get
%	\begin{equation}\label{eq:Subtraction}
%		\begin{aligned}
%			0 &\geq  - \int_{\Omega}|\nabla(u_{\lambda'}-w_{\lambda})|_{\mathbb{G}}^2 \, \theta'_{\varepsilon}(w_{\lambda}-u_{\lambda'})\\
%			&= \int_{\Omega}\left( \dfrac{\lambda'}{u_{\lambda'}^{\gamma}} - \dfrac{\lambda}{w_{\lambda}^{\gamma}} + u_{\lambda'}^{2^{*}-1}\right) \theta_{\varepsilon}(w_{\lambda}-u_{\lambda'})\\
%			&\geq \lambda \int_{\Omega}\left( \dfrac{1}{u_{\lambda'}^{\gamma}}- \dfrac{1}{w_{\lambda}^{\gamma}}\right)\theta_{\varepsilon}(w_{\lambda}-u_{\lambda'}).
%		\end{aligned}
%	\end{equation}
%	Now, letting $\varepsilon \to 0^+$ we find that
%	$$\int_{\{w_{\lambda}> u_{\lambda'}\}}\left( \dfrac{1}{u_{\lambda'}^{\gamma}}- \dfrac{1}{w_{\lambda}^{\gamma}}\right)\leq 0,$$
%	\noindent and this implies that
%	$$\left|\left\{ g \in \Omega: w_{\lambda}(g) > u_{\lambda'}(g)\right\}\right| = 0,$$
%	\noindent as claimed in \eqref{eq:Claim}. \vspace*{0.1cm}
	
We now set $\overline{u}= u_{\lambda'}$ and $\underline{u} = w_{\lambda}$,	
	and we apply Lemma \ref{lem:Perron}: this immediately yields that problem \eqref{EqProblem} admits a weak solution $u_{\lambda}$ for every $\lambda \in (0,\Lambda)$. Moreover, recalling the definition of $I_{\lambda}$ in \eqref{eq:Def_Ilambda}, such a solution satisfies that
	$$I_{\lambda}(u_\lambda) = \min\{u\in S_0^1(\Omega):\,w_{\lambda}\leq u\leq u_{\lambda'}\}
	\leq I_{\lambda}(w_{\lambda}).$$
	In particular,
	by Theorem \ref{thm:Sublinear_Problem} we have
	\begin{equation} \label{eq:Ilambdaulambdaneg}
		I_{\lambda}(u_{\lambda}) \leq I_{\lambda}(w_{\lambda}) \leq J_{\lambda}(w_{\lambda}) <0.
	\end{equation}
	It remains to consider the case $\lambda = \Lambda$. The proof
	is rather standard and pretty similar to that of \cite[Lemma 3.5]{BiGaVe}. We report it here for the sake of completeness.
	To begin with, we choose a monotone increasing sequence $\{\lambda_k\}_k\subseteq(0,\Lambda)$ such that 
     $\lambda_{k} \to \Lambda$ as $k\to+\infty$. Now, for each $k \in \mathbb{N}$, we set
	$$u_k := u_{\lambda_k}\in S^{1}_{0}(\Omega),$$
	\noindent where $u_{\lambda_k}$ is the weak solution of problem \eqref{EqProblem} (with $\lambda = \lambda_k$) constructed
	as above by means of Lemma \ref{lem:Perron}. Thanks to 
	\eqref{eq:Ilambdaulambdaneg}, for every $k\geq 1$ we have
	\begin{equation} \label{eq:Ilambdakneg}
		I_{\lambda_k}(u_{k}) 
		= \dfrac{1}{2} \int_{\Omega}|\nabla_{\mathbb{G}}u_k|^2 - \dfrac{\lambda_k}{q+1}\int_{\Omega}|u_k|^{q+1} - \dfrac{1}{2^{\star}_{Q}}\int_{\Omega}|u_k|^{2^{\star}_{Q}} <0.
	\end{equation}
	Moreover, by using $\varphi = u_k$ in \eqref{eq:weak_sub_super_sol},  and recalling that $u_k$
	solves \eqref{EqProblem} (with $\lambda = \lambda_k$), we get
	\begin{equation} \label{eq:testwithukzero}
		\int_{\Omega}|\nabla_{\mathbb{G}}u_k|^2 -\lambda_k\int_\Omega u_k^{q+1}
		-\int_\Omega u_k^{2^{\star}_{Q}} = 0.
	\end{equation}
	Combining \eqref{eq:Ilambdakneg} with \eqref{eq:testwithukzero},
	we notice that the sequence $\{u_k\}_k$ is  bounded in~$S^{1}_{0}(\Omega)$.
	Therefore, we can find a function
	$$u_{\Lambda}\in S^{1}_{0}(\Omega),$$ 
	such that
	(up to a subsequence and as $k\to+\infty$)
	\begin{itemize}
		\item[a)] $u_k\to u_{\Lambda}$ weakly in $S^{1}_{0}(\Omega)$ and strongly
		in $L^p(\Omega)$ for $1\leq p <2^{\star}_{Q}$;
		\item[b)] $u_k\to u_{\Lambda}$ a.e.\,in $\Omega$.
	\end{itemize}
	We now observe that, being $\{\lambda_k\}_k$ increasing, it follows that $\lambda_k\geq \lambda_1$ for every $k\geq 1$.
	Moreover, arguing as above yields that
	$u_{\lambda_k}\geq w_{\lambda_1}$, and thus
	$$u_{\Lambda} > 0\quad\text{a.e.\,in $\Omega$}.$$
	Moreover, since $u_k$ solves problem \eqref{EqProblem} (with $\lambda = \lambda_k$),
	we have
	$$\int_{\Omega}\langle \nabla_{\mathbb{G}}u_k,\nabla_{\mathbb{G}}\varphi\rangle_{\mathfrak{g}_{1}} -\lambda_k\int_\Omega u_k^{q}\varphi 
	-\int_\Omega u_k^{2^{\star}_{Q}-1}\varphi = 0\quad\text{for every $\varphi\in S^{1}_{0}(\Omega)$}.$$
Therefore, passing to the limit as $k\to+\infty$ in the above identity, and by dominated convergence, we get that $u_{\Lambda}$ satisfies
	$$\int_{\Omega}\langle \nabla_{\mathbb{G}}u_{\Lambda},\nabla_{\mathbb{G}}\varphi\rangle_{\mathfrak{g}_{1}} -\Lambda\int_\Omega u_{\Lambda}^{q}\varphi
	-\int_\Omega u_{\Lambda}^{2^{\star}_{Q}-1}\varphi = 0\quad\text{for every $\varphi\in S^{1}_{0}(\Omega)$},$$
	\noindent which shows that $u_{\Lambda}$ is actually a weak solution of
	problem \eqref{EqProblem} (with $\lambda = \Lambda$). 
%	
%	We explicitly
%	point out that the convergence of the 
%	$\int_{\Omega}\langle \nabla_{\mathbb{G}}u_k,\nabla_{\mathbb{G}}\varphi\rangle_{\mathfrak{g}_{1}}$ 
%	to $\int_{\Omega}\langle \nabla_{\mathbb{G}}u_{\Lambda},\nabla_{\mathbb{G}}\varphi\rangle_{\mathfrak{g}_{1}}$ follows
%	from the weak convergence of $u_k$ to $u_{\Lambda}$, since 
%	$$v\mapsto \int_{\Omega}\langle \nabla_{\mathbb{G}}v,\nabla_{\mathbb{G}}\varphi\rangle_{\mathfrak{g}_{1}}$$
%	is a linear and continuous functional on $S^{1}_{0}(\Omega)$. 
	This closes the proof.
\end{proof}

\medskip

Before tackling the problem of a second solution, we focus
on the behavior of the functional $I_\lambda$ around $u_\lambda$. In particular we will show that the first solution $u_{\lambda}$ is actually a local minimum in the $S_{0}^{1}(\Omega)$-topology. As recalled in the Introduction, in the Euclidean case this is performed exploiting a famous result by Brezis and Nirenberg \cite{BNH1C1}. In our case we have to follow a different approach due to the lack of boundary regularity of the solution. In particular, we adapt the strategy used in \cite{AbDiVa} which in turn is inspired by a work of Alama.

\begin{lemma}\label{lem:locmin}
	For $\lambda\in (0,\Lambda)$, if $u_\lambda$ is the solution
	presented in Theorem \ref{thm:Existence_First}, then
	$u_\lambda$ is a local minimum for $I_\lambda$ in the $S^1_0(\Omega)$-topology,
	meaning that there exists $r_0>0$ such that for any $u\in S^1_0(\Omega)$,
	\[
	I_\lambda(u_\lambda)\leq I_\lambda(u) \quad \textrm{ for all } u\in S^{1}_{0}(\Omega) \textrm{ with } \|u-u_\lambda\|_{S_{0}^{1}(\Omega)}< r_0 .
	\]
\end{lemma}
\begin{proof}

	In the following of this proof, we denote by $\overline u$ the solution
	of \eqref{EqProblem} (with $\lambda = \overline \lambda$) for some value $\overline \lambda$
	such that  $\lambda<\overline \lambda <\Lambda$.
	
%	We introduce the convex and closed subset $M:=\{u\in S^1_0(\Omega):\ 0\leq u\leq \overline u\}\subset S^1_0(\Omega)$
%	and consider the infimum of $I_\lambda$ over it, in particular there exists
%	$u_*\in M$ such that
%	\[
%	I_\lambda(u_*)=\inf_{u\in M}I_\lambda(u).
%	\]
%	
%	Observe that $u_*\neq0$. Indeed, if $v_*$ is the unique solution of
%		\[
%	\left\{
%	\begin{array}{rll}
%		-\Delta_\G v&=\lambda v^q&\m{in }\Omega\\
%		v&\geq0 &\m{in }\Omega\\
%		v&=0 &\m{on }\de\Omega.
%	\end{array}
%	\right.
%	\]
%	then $I_\lambda(v_*)<0$ and so $u_*$ must be non-null.
%	
%	The result in STRUWE p.17 tells us that $u_*$ is a solution to \eqref{EqProblem}$_\lambda$.
%	If $u_*\neq u_\lambda$ we just found a second solution to \eqref{EqProblem}$_\lambda$, thus proving
%	THEOREM REF for $\lambda$. Thereby we can address the case $u_*=u_\lambda$, meaning
%	that the minimal solution of \eqref{EqProblem}$_\lambda$ obtained
%	via Lemma \ref{lem:minsol}
%	is a constrained minimum for $I_\lambda$ on $M$.\newline
	
	Let's suppose by contradiction that there exists a sequence $\{v_n\}$ in $S^1_0(\Omega)$ verifying that
	$\|v_n-u_\lambda\|_{S_{0}^{1}(\Omega)}\to0$ and $I_\lambda(v_n)<I_\lambda(u_\lambda)$ for each $n$.
	We also introduce two ausiliary functions,
	\begin{align*}
		w_n&:=(v_n-\overline u)_+\, ,\\
		u_n&:=\max(0,\min(v_n,\overline u)),
	\end{align*}
	and the sets 
	\begin{align*}
		T_n&:=\{x\in \Omega:\ u_n(x)=v_n(x)\},\\
		S_n&:=\supp(w_n)\cap\Omega.
	\end{align*}
	
	We want to prove that 
	\begin{equation}\label{Sn0}
		\lim_{n\to+\infty}|S_n|=0.
	\end{equation}
	Consider the two sets,
	\begin{align*}
		E(n,\delta)&:=\{x\in \Omega:\ v_n(x)\geq\overline u(x)>u_\lambda(x)+\delta\},\\
		F(n,\delta)&:=\{x\in \Omega:\ v_n(x)\geq\overline u(x),\ \m{and}\ \overline u(x)\leq u_\lambda(x)+\delta\}.
	\end{align*}
	By construction, $S_n\subset E_n(n,\delta)\cup F(n,\delta)$ for each $n$ and each $\delta>0$.
	We are going to show that
	for any $\varepsilon>0$ and a suitable choice of $\delta>0$ and $n\in \N$, we have both $|E(n,\delta)|,|F(n,\delta)|<\frac{\varepsilon}{2}$.
	\begin{enumerate}
		\item We start from $E(n,\delta)$.
		By definition of the sequence $v_n$, we have $\|v_n-u_\lambda\|_{L^2(\Omega)}\to0$.
		Therefore if we fix $\delta$ and $\varepsilon$ there exists $n_0$ such that for any $n\geq n_0$,
		$\frac{\delta^2\varepsilon}{2}>\|v_n-u_\lambda\|_{L^2}$. As a consequence,
		\[
		\frac{\delta^2\varepsilon}{2}>\int_\Omega |v_n-u_\lambda|^2\geq\int_{E(n,\delta)}|v_n-u_\lambda|^2>\delta^2\cdot|E(n,\delta)|.
		\]
		This implies that $|E(n,\delta)|<\frac{\varepsilon}{2}$ for any $\delta$ and any $n\geq n_0$.
		
		\item Let's consider $F(n,\delta)$. If,
		\begin{align*}
			\overline F(\delta)&:=\{x\in \Omega:\ \overline u(x)\leq u_\lambda(x)+\delta\}\\
			\overline F&:=\{x\in \Omega:\ \overline u(x)\leq u_\lambda(x)\},
		\end{align*}
		then by construction
		\[
		0=|\overline F|=\left|\bigcap_{m=1}^\infty \overline F\left(\frac{1}{m}\right)\right|=\lim_{m\to+\infty}\left|\overline F\left(\frac{1}{m}\right)\right|.
		\]
		Therefore, for a suitable $m_0$ we get that $|\overline F(\delta)|<\frac{\varepsilon}{2}$ for any $\delta<\frac{1}{m_0}$
		and \cor{a fortiori} $|F(n,\delta)|<\frac{\varepsilon}{2}$ for any $n$ because $F(n,\delta)\subset\overline F(\delta)$.\newline
	\end{enumerate}
	
	This proves \eqref{Sn0}. Let's now consider the function
	\[
	h(u):=\frac{\lambda}{q+1}u_+^{q+1}+\frac{u_+^{2^\star_Q}}{2^\star_Q},
	\]
	and observe that by definition $\Omega=T_n\cup S_n$ because $u_n\leq v_n$.
	We develop the evaluations,
	\begin{align*}
		I_\lambda(v_n)&=\frac{1}{2}\|v_n\|_{S_{0}^{1}(\Omega)}^2-\int_\Omega h(v_n)\\
		&\geq \frac{1}{2}\|v_n^+\|_{S_{0}^{1}(\Omega)}^2+\frac{1}{2}\|v_n^-\|_{S_{0}^{1}(\Omega)}^2-\int_{T_n} h(v_n)-\int_{S_n}h(v_n)\\
		&(\m{because }u_n=\overline u\m{ and }v_n=w_n+\overline u\m{ on }S_n)\\
		&=\frac{1}{2}\|v_n^+\|_{S_{0}^{1}(\Omega)}^2+\frac{1}{2}\|v_n^-\|_{S_{0}^{1}(\Omega)}^2-\int_\Omega h(u_n)-\int_{S_n}(h(w_n+\overline u)-h(\overline u))\\
		&=I_\lambda(u_n)+\frac{1}{2}\left(\|v_n^+\|_{S_{0}^{1}(\Omega)}^2-\|u_n\|_{S_{0}^{1}(\Omega)}^2\right)+\frac{1}{2}\|v_n^-\|_{S_{0}^{1}(\Omega)}^2-\int_{S_n}\left(h(w_n+\overline u)-h(\overline u)\right).
	\end{align*}
	Here as $v_n^+=u_n+w_m$, using that $I_\lambda(u_n)\geq I_\lambda (u_\lambda)$
	because $u_n\in M$ for any $n$, and that $\overline u$ is a supersolution of \eqref{EqProblem},
	we get 
	\begin{equation}\label{eqADV}
		\begin{aligned}
		I_\lambda(v_n)&\geq I_\lambda (u_\lambda)+\frac{1}{2}\|w_n\|_{S_{0}^{1}(\Omega)}^2+\langle u_n,w_n\rangle_{S_{0}^{1}(\Omega)}+\frac{1}{2}\|v_n^-\|_{S_{0}^{1}(\Omega)}^2
		-\int_{S_n}\left(h(w_n+\overline u)-h(\overline u)\right)\\
		&\geq I_\lambda(u_\lambda)+\frac{1}{2}\|w_n\|_{S_{0}^{1}(\Omega)}^2+\frac{1}{2}\|v_n^-\|_{S_{0}^{1}(\Omega)}^2
		-\int_{S_n}\left(h(w_n+\overline u)-h(\overline u)-\lambda \overline u^qw_n-\overline u^{2^\star_Q-1}w_m\right).
		\end{aligned}
	\end{equation}
	
	Let's use the notation
	\[
	f_r(x,y)=\frac{1}{r+1}(x+y)^{r+1}-\frac{1}{r+1}y^{r+1}-xy^r.
	\]
	As $w_n\geq0$ and $\overline u\geq0$, simply
	by Taylor expasion we have
	\begin{equation}\label{tayq}
		0\leq f_{q}(w_n,\overline u)\leq \frac{q}{2}w_n^2\overline u^{q-1}.
	\end{equation}
		Mooreover, by \cite[Theorem 3.4]{Ruzhansky_Suragan}, we have 
	\begin{equation}\label{picone}
		 \lambda\int_\Omega \overline u^{q-1}w_n^2\leq \int_\Omega (-\Delta \overline u)\frac{w_n^2}{\overline u}\leq \|w_n\|_{S_{0}^{1}(\Omega)}^2,
	\end{equation}
	and combining the last two equations, we get
	\begin{equation}\label{combq}
		\lambda\int_\Omega f_q(w_n,\overline u)\leq \frac{q}{2}\int_\Omega w_n^2\overline u^{q-1}\leq \frac{q}{2}\|w_n\|_{S_{0}^{1}(\Omega)}^2.
	\end{equation}
%	\begin{equation}\label{tayq}
%		0\leq \frac{1}{q+1}(w_n+\overline u)^{q+1}-\frac{1}{q+1}\overline u^{q+1}-\overline u^qw_n\leq \frac{q}{2}\overline u^{q-1}w_n^2.
%	\end{equation}
 	By similar reasoning and using Sobolev inequality, we obtain
 \begin{equation}
 	\int_\Omega f_{2^\star_Q-1}(w_n,\overline u)\leq o(1)\cdot \|w_n\|_{S_{0}^{1}(\Omega)}^2.
 \end{equation}
 As $h(w_n+\overline u)=f_q(w_n,\overline u)+f_{2^\star_Q-1}(w_n,\overline u)$, 
 by combining the last results with \eqref{eqADV},
 we obtain
 \[
 I_\lambda(v_n)\leq I_\lambda(u_\lambda)+\frac{1}{2}\|w_n\|_{S_{0}^{1}(\Omega)}^2(1-q-o(1))+\frac{1}{2}\|v_n^-\|_{S_{0}^{1}(\Omega)}.
 \]
 As $I_\lambda(v_n)<I_\lambda(u_\lambda)$ by hypothesis,
 and $q<1$, then for $n$ sufficiently big we must have $v_n^-=0$,
 but this imply $v_n\geq0$ and therefore $v_n\in M$,
 thus contradicting $I_\lambda(v_n)<I_\lambda(u_\lambda)$.
 This completes the proof.	\end{proof}
\medskip

\color{black}



\medskip


\section{Proof of Theorem \ref{thm:main} - Part C)}\label{sec:Second_Solution}

The aim of this section is to prove that the problem \eqref{EqProblem} actually admits a second
solution for every $\lambda\in (0,\Lambda)$.
We briefly recall that, thanks to Theorem \ref{thm:Existence_First} and Lemma \ref{lem:locmin}, we have that 
\begin{itemize}
	\item for every $\lambda \in (0,\Lambda]$, the problem \eqref{EqProblem} admits a solution denoted by $u_{\lambda}$;
	\item for every $\lambda \in (0,\Lambda)$, $u_{\lambda}$ is a local minimum for $I_{\lambda}$ in the $S_{0}^{1}$-topology, i.e.
	there exists $r_0>0$ such that for any $u\in S^1_0(\Omega)$,
	\[
	I_\lambda(u_\lambda)\leq I_\lambda(u) \quad \textrm{ for all } u\in S^{1}_{0}(\Omega) \textrm{ with } \|u-u_\lambda\|_{S_{0}^{1}(\Omega)}< r_0 .
	\]
\end{itemize}

We now adapt to our setting the strategy used in \cite{Tarantello}.
Firstly, we are going to consider two cases:
\begin{enumerate}
	\item for every $r\in (0,r_0)$,
	\[
	\inf_{\|u-u_\lambda\|_{S_{0}^{1}(\Omega)}=r}I_\lambda(u)=I_\lambda(u_\lambda),
	\]
	\item there exists $r\in (0,r_0)$ such that 
	\[
	\inf_{\|u-u_\lambda\|_{S_{0}^{1}(\Omega)}=r}I_\lambda(u)>I_\lambda(u_\lambda).
	\]
\end{enumerate}
We treat separately the two cases in the following sub-sections. 
%for the sake of a clearer notation, we set $p:=2^*_Q-1$.
\newline

\subsection{First case}\label{subsec:CaseI}
%We impose that $\inf I_\lambda(u)=I_\lambda(u_\lambda)$
%for every $r\in (0,r_0)$, where
%we are setting $\|u-u_\lambda\|_{S_{0}^{1}(\Omega)}=r$.

%We start by observing that for $p\in (1,2^*_Q-1)$
%the functional $I_\lambda$ respects the Palais-Smale condition (REFERENCE?),
%therefore we know (DEFIGUEREIDO) that 
%for every $r\in (0,r_0)$ the
%infimum is always reached at a critical point $w_\lambda$ of~$I_\lambda$
%such that $\|w_\lambda-u_\lambda\|=r$.
%Thus in this case we have a second distinct solution $w_\lambda$ of \eqref{EqProblem}$_\lambda$.\newline

%At the critical
%value, $I_\lambda$ lacks the PS condition, therefore we apply a method very similar (REFERENCE)
%to recover a second solution. 
We introduce the space
\begin{equation}\label{def:hlambda}
	H_\lambda:=\{u\in S^1_0(\Omega):\,u\geq u_\lambda \m{ a.e.~in }\Omega\}.
\end{equation}

\noindent By the standing assumptions in (1), there exists 
$\{u_k\}_k\subset H_{\lambda}$ such that 
\begin{enumerate}[i)]
	\item $\|u_k-u_\lambda\|_{S^1_0(\Omega)} = r$ for every $k\geq 1$;	
	\item as $k\to+\infty$ we have $I_\lambda(u_k)\to I_\lambda(u_\lambda)$.
\end{enumerate}
We also introduce the space
$$X_\lambda := \{u\in H_\lambda:\,r-\bar r\leq\|u-u_\lambda\|_{S^1_0(\Omega)}\leq r+\bar r\},$$
where $\bar r>0$ is taken sufficiently small to have $r-\bar r>0$ and $r+\bar r<r_0$.
The set $X_\lambda$ becomes a complete metric space once endowed with the distance associated to the norm $\|\cdot\|_{S^1_0(\Omega)}$.

We can now proceed very similarly to \cite{BiGaVe},
by applying the Ekeland's Variational Principle 
and therefore obtaining a sequence
 $\{w_k\}_k\subset X_\lambda$ such that
\begin{equation} \label{eq:EkelandCaseA}
	\begin{split}
		\mathrm{i)}&\,\, I_\lambda(w_k)\leq I_\lambda(u_k)\leq I_\lambda(u_\lambda)+\frac{1}{k^2}, \\
		\mathrm{ii)}&\,\,\|w_k-u_k\|_{S^1_0(\Omega)}\leq \frac{1}{k}, \\
		\mathrm{iii)}&\,\, I_\lambda(w_k)\leq I_\lambda(u)+\frac{1}{k}\,
		\|w_k-u\|_{S^1_0(\Omega)}\quad\text{for every $u\in X_\lambda$}.    
	\end{split}
\end{equation}
By boundedness of $\{w_k\}$ in $S^1_0(\Omega)$, there exists
$w_\lambda\in S^1_0(\Omega)$ such that 
the following are true (up to a sub-sequence),
\begin{equation} \label{eq:limitvkCaseA}
	\begin{split}
		\mathrm{i)}&\,\,\text{$w_k\to w_\lambda$ weakly in $S^1_0(\Omega)$}; \\
		\mathrm{ii)}&\,\,\text{$w_k\to w_\lambda$ strongly in $L^r(\Omega)$ for every $1\leq r<2^\star_Q$}; \\
		\mathrm{iii)}&\,\,\text{$w_k\to w_\lambda$ pointwise
			a.e.\,in $\Omega$}.
	\end{split}
\end{equation}
%where we have also used the \emph{compact embedding} 
%$S^1_0(\Omega)\hookrightarrow L^2(\Omega)$ (see Section \ref{sec:fsandpdes}).

First we show that the limit function $w_{\lambda}$ is a solution to \eqref{EqProblem}. This is the content of the following

\begin{lemma}\label{lem:casoa1}
	The function $w_\lambda$ is a weak solution of \eqref{EqProblem}.
\end{lemma}
\proof
Given $w\in H_\lambda$ we consider $\varepsilon_0$
sufficiently small that $w_k+\varepsilon(w-w_k)\in X_\lambda$ for each $0<\varepsilon<\varepsilon_0$.
For $k$ sufficiently big such an $\varepsilon_0$ always exists 
because, indeed
\begin{multline*}
 r-\frac{1}{k}\leq \|u_k-u_\lambda\|_{S^1_0(\Omega)}-\|w_k-u_k\|_{S^1_0(\Omega)}\leq \|w_k-u_\lambda\|_{S^1_0(\Omega)}\\
 \leq \|u_k-u_\lambda\|_{S^1_0(\Omega)}+\|w_k-u_k\|_{S^1_0(\Omega)}\leq r+\frac{1}{k}.
\end{multline*}
By setting $u=w_k+\varepsilon(w-w_k)$ in \eqref{eq:EkelandCaseA}
we get
\[
\frac{I_\lambda(w_k+\varepsilon(w-w_k))-I_\lambda(w_k)}{\varepsilon}\geq -\frac{1}{k}\, \|w-w_k\|_{S^1_0(\Omega)}.
\]
Letting $\varepsilon\to0^+$, we get
\begin{multline}\label{eq:ineq2.21haitao}
		 -\frac{1}{k}\, \|w-w_k\|_{S^1_0(\Omega)}\leq \int_\Omega\langle \nabla_\G w_k,\nabla_\G(w-w_k)\rangle_{\mathfrak{g}_{1}}\\
	-\int_\Omega w_k^{2^\star_Q-1}(w-w_k)
		-\lambda \int_\Omega w_k^q(w-w_k)\quad\textrm{ for every } w\,\,\in H_\lambda.
\end{multline}
%with $0<\theta<1$. Observe that $v_k+\varepsilon(w-v_k)\geq u_\lambda$ a.e.~in $\Omega$,
%and $w-v_k\in S^1_0(\Omega)$, 
%therefore 
%\[
%\int_\Omega \left| (v_k+\theta\varepsilon(w-v_k))^{-\gamma}(w-v_k)\right|\leq \int_\Omega u_\lambda^{-\gamma}|w-v_k|<\infty.
%\]
%By the dominated convergence theorem, this implies
%\[
%\lim_{\varepsilon\to0^+}\int_\Omega(v_k+\theta\varepsilon(w-v_k))^{-\gamma}(w-v_k)=\int_\Omega v_k^{-\gamma}(w-v_k),
%\]
%and therefore we get the inequality
%\begin{equation}\label{eq:ineq2.21haitao}
%\begin{split}
%	& -\frac{1}{k}\, \|w-v_k\|_{S^1_0(\Omega)}\\
%	& \qquad\leq \int_\Omega\langle \nabla_\G v_k,\nabla_\G(w-v_k)\rangle_{\mathfrak{g}_{1}}-\int_\Omega v_k^{2^\star_Q-1}(w-v_k)
%	-\lambda \int_\Omega v_k^{-\gamma}(w-v_k)\ \ \forall w\,\,\in H_\lambda.
%	\end{split}
%	\end{equation}
For any $\varphi\in S^1_0(\Omega)$ and $\varepsilon>0$ we introduce the functions
$$\varphi_{k,\varepsilon}:=w_k+\varepsilon\varphi-u_\lambda \quad \textrm{ and } \quad 
\varphi_{\varepsilon}:=w_\lambda+\varepsilon\varphi-u_\lambda.$$
We further set $w:=w_k+\varepsilon\varphi+(\varphi_{k,\varepsilon})_-\in H_\lambda$, 
then, by \eqref{eq:ineq2.21haitao},
\begin{multline}\label{eq:ineq2.22haitao}
-\frac{1}{k}\, \|\varepsilon\varphi+(\varphi_{k,\varepsilon})_-\|_{S^1_0(\Omega)}\leq \int_\Omega\langle \nabla_\G w_k,\nabla_\G(\varepsilon\varphi+(\varphi_{k,\varepsilon})_-)\rangle_{\mathfrak{g}_{1}}\\
-\int_\Omega w_k^{2^\star_Q-1}(\varepsilon\varphi+(\varphi_{k,\varepsilon})_-)
-\lambda\int_\Omega w_k^{q}(\varepsilon\varphi+(\varphi_{k,\varepsilon})_-).
\end{multline}
From \eqref{eq:limitvkCaseA},
$$(\varphi_{\varepsilon,k})_-\to(\varphi_\varepsilon)_-  \quad \textrm{ a.e. in } \Omega, \textrm{ as } k\to+\infty.$$
Moreover, we have
%\begin{itemize}
%	\vspace{7pt}
%	\item $v_k^{2^\star_Q-1}\, (\varphi_{k,\varepsilon})_-=v_k^{2^\star_Q-1}\, (u_\lambda-\varepsilon\varphi-v_k)\cdot \mathbf{1}_{\{u_\lambda-\varepsilon\varphi-v_k\geq0\}}
%		\leq(u_\lambda+\varepsilon|\varphi|)^{2^\star_Q}$
%		\vspace{7pt}
%		\item $v_k^{-\gamma}\, (\varphi_{k,\varepsilon})_-\leq 2\varepsilon \, |\varphi|\, u_\lambda^{-\gamma}$.
%		\end{itemize}
\[
w_k^{2^\star_Q-1}(\varphi_{k,\varepsilon})_-=w_k^{2^\star_Q-1}(u_\lambda-\varepsilon\varphi-w_k)\cdot \mathbf{1}_{\{u_\lambda-\varepsilon\varphi-w_k\}}
\leq (u_\lambda+\varepsilon|\varphi|)^{2^\star_Q}.
\]
%By the results above and Remark	\ref{rem:defweaksolPb}-(2), we can use the Dominated
%Convergence Theorem, obtaining
Therefore, by Dominated Convergence we get
\begin{equation}\label{eq:gather1}
	\begin{split}
	\lim_{k\to+\infty}\biggl(\int_\Omega w_k^{2^\star_Q-1}(\varepsilon\varphi+(\varphi_{k,\varepsilon})_-)
	&+\lambda\int_\Omega w_k^{q}(\varepsilon\varphi+(\varphi_{k,\varepsilon})_-)\biggl)\\
	&=\int_\Omega w_\lambda^{2^\star_Q-1}(\varepsilon\varphi+(\varphi_\varepsilon)_-)+
	\lambda\int_\Omega w_\lambda^{q}(\varepsilon\varphi+(\varphi_\varepsilon)_-)
	\end{split}
\end{equation}
For the other term, similarly to  \cite[Lemma 3.4]{BadTar}
we have
\[
\int_\Omega \langle \nabla_\G w_k,\nabla_\G (\varphi_{k,\varepsilon})_-\rangle_{\mathfrak{g}_{1}}\leq
\int_\Omega\langle \nabla_\G{w_\lambda},\nabla_\G(\varphi_{\varepsilon})_-\rangle_{\mathfrak{g}_{1}} +o(1)\ \m{as}\ k\to+\infty
\]
and because $v_k\rightharpoonup v_\lambda$  weakly in $S_0^1(\Omega)$, we obtain
\begin{equation}\label{eq:gather2}
	\int_\Omega\langle \nabla_\G w_k,\nabla_\G(\varepsilon\varphi+(\varphi_{k,\varepsilon})_-)\rangle_{\mathfrak{g}_{1}}
	\leq \int_\Omega\langle \nabla_\G {w_\lambda},\nabla_\G(\varepsilon\varphi+(\varphi_{\varepsilon})_-)\rangle_{\mathfrak{g}_{1}} +o(1)\ \m{as}\ k\to+\infty.
\end{equation}
As $\|w_k\|_{S^1_0(\Omega)}$ is uniformly bounded
w.r.t. $k$, we have the same
for $\|(\varphi_{k,\varepsilon})_-\|_{S^1_0(\Omega)}$.
Therefore we can pass to the limit as $k\to+\infty$
 in \eqref{eq:ineq2.22haitao}, 
and by recalling \eqref{eq:gather1} and \eqref{eq:gather2}  we obtain
\begin{equation}
	\int_\Omega\langle \nabla_\G w_\lambda,\nabla_\G(\varepsilon\varphi+(\varphi_\varepsilon)_-)\rangle_{\mathfrak{g}_1}\geq
	\int_\Omega w_\lambda^{2^\star_Q-1}(\varepsilon\varphi+(\varphi_\varepsilon)_-)+
	\lambda\int_\Omega w_\lambda^{q}(\varepsilon\varphi+(\varphi_\varepsilon)_-).
\end{equation}
Finally, we can conlude as in \cite[Lemma 4.1]{BiGaVe}, getting
%Exploiting the computations carried out in \cite[Lemma 2.6]{Haitao}, we get
%\[
%\begin{split}
%	\int_\Omega\langle \nabla_\G v_\lambda,\nabla_\G \varphi\rangle_{\mathfrak{g}_{1}} 
%	&- \lambda \int_\Omega v_\lambda^{-\gamma}\varphi-\int_\Omega v_\lambda^{2^\star_Q-1}\varphi\\
%	&\geq-\frac{1}{\varepsilon} \left(\int_\Omega \langle \nabla_\G v_\lambda,\nabla_\G (\varphi_\varepsilon)_-\rangle_{\mathfrak{g}_{1}}
%	-\lambda\int_\Omega v_\lambda^{-\gamma}(\varphi_\varepsilon)_--\int_\Omega v_\lambda^{2^\star_Q-1}(\varphi_\varepsilon)_-\right)\\
%	&(\m{since }u_\lambda\m{ is a solution of \eqref{EqProblem}}_\lambda)\\
%	&=-\frac{1}{\varepsilon} \biggl(\int_\Omega \langle \nabla_\G(v_\lambda-u_\lambda),\nabla_\G(\varphi_\varepsilon)_-\rangle_{\mathfrak{g}_{1}}
%	-\lambda\int_\Omega(v_\lambda^{-\gamma}-u_\lambda^{-\gamma})(\varphi_\varepsilon)_-\\
%	&\ \ \ 	-\int_\Omega(v_\lambda^{2^\star_Q-1}-u_\lambda^{2^\star_Q-1})(\varphi_\varepsilon)_-\biggl)\\
%	\vspace{6pt}
%	&(\m{since }v_\lambda\geq u_\lambda\m{ a.e.~and }v_\lambda=\lim_{k\to+\infty}v_k)\\
%	&\geq -\frac{1}{\varepsilon} 
%	\biggl(\int_{\{u_\lambda> v_\lambda+\varepsilon\varphi\}}\langle \nabla_\G(v_\lambda-u_\lambda),\nabla_\G (v_\lambda-u_\lambda+\varepsilon\varphi)\rangle_{\mathfrak{g}_{1}}\\
%	&\ \ \ 	-\lambda\int_{\{u_\lambda> v_\lambda+\varepsilon\varphi\}}(v_\lambda^{-\gamma}-u_\lambda^{-\gamma})(v_\lambda-u_\lambda+\varepsilon\varphi)\biggl)\\
%	&\geq o(1)\ \ \ \m{as}\ \varepsilon\to0^+,
%\end{split}
%\]
%\noindent where we used \eqref{eq:Grad_and_level_set}
%and, as before, that
%\[
%\bigcap_{\varepsilon>0}\{u_\lambda> v_\lambda+\varepsilon\varphi\}=\varnothing.
%\]
%This implies that $|\{u_\lambda>v_\lambda+\varepsilon\varphi\}|\to 0$
%and therefore, by letting $\varepsilon\to0^+$ in the inequality above, we conclude that
\begin{equation}
	\int_\Omega\langle \nabla_\G w_\lambda,\nabla_\G\varphi\rangle_{\mathfrak{g}_{1}}-\lambda\int_\Omega w_\lambda^{q}\varphi
	-\int_\Omega v_\lambda^{2^\star_Q-1}\varphi\geq0,
\end{equation}
and by the arbitrariness of $\varphi \in S^1_0(\Omega)$ 
we conclude that $w_\lambda$
is a weak solution of \eqref{EqProblem}.\fine

\color{black}

%It remains to show that $v_\lambda$ and $u_\lambda$ are different solutions.
%In order to do that we use the following lemma.
The following technical Lemma closely follows \cite[Lemma 4.2]{BiGaVe}.

\begin{lemma}\label{lem:casoa2}
	If $w_\lambda$ is as above,
	then $\|w_\lambda-u_\lambda\|_{S^{1}_{0}(\Omega)}=r$.
\end{lemma}
\begin{proof} 
We notice first that it is enough to prove that 
\begin{equation}\label{eq:stronconv}
	w_k\to w_\lambda\ \m{strongly in }S_0^1(\Omega)\ \m{as}\ k\to+\infty.
\end{equation}
Indeed, by using that $\|u_k-u_\lambda\|_{S^{1}_{0}(\Omega)}=r$ for any $k$,
we get
\[
r-\|w_k-u_k\|_{S^{1}_{0}(\Omega)} \leq \|w_k-u_\lambda\|_{S^{1}_{0}(\Omega)} \leq r+\|w_k-u_k\|_{S^{1}_{0}(\Omega)}.
\]
Combining the latter with both the strong convergence in $S_{0}^{1}(\Omega)$ and \eqref{eq:EkelandCaseA}-ii),
%the fact that $\|w_k-u_k\|_{S^{1}_{0}(\Omega)}\leq \frac{1}{k}$
implies 
$$\|w_\lambda-u_\lambda\|_{S^{1}_{0}(\Omega)}=r.$$
Let us now proceed with the proof of \eqref{eq:stronconv}. 
In view of \eqref{eq:limitvkCaseA}, and arguing as in the proof of \cite[Lemma 4.3]{BiGaVe}, it follows that


%%Since $w_k\to w_\lambda$ weakly in $S_0^1(\Omega)$ as $k\to+\infty$, we get that
%
%Moreover, arguing exactly as in the proof of \cite[Lemma 4.3]{BiGaVe}, we also get
\begin{align}
	\|w_k - w_\lambda\|_{L^{q+1}(\Omega)} &\to 0 \quad \textrm{ as } k \to +\infty, \label{eq:case1concl1}\\
	\|w_k\|^{2^\star_Q}_{L^{2^\star_Q}}(\Omega)&=\|w_\lambda\|_{L^{2^\star_Q}(\Omega)}^{2^\star_Q}+
	\|w_k-w_\lambda\|_{L^{2^\star_Q}(\Omega)}^{2^\star_Q}+o(1),\\
	\|w_k\|_{S^{1}_{0}(\Omega)}^2&=\|w_\lambda\|_{S^{1}_{0}(\Omega)}^2+\|w_k-w_\lambda\|_{S^{1}_{0}(\Omega)}^2+o(1).\label{eq:case1concl2}
\end{align}
In particular, from \eqref{eq:case1concl1} we get
\begin{equation}\label{eq:case1concl1bis}
	\int_\Omega w_k^{q+1}=\int_\Omega w_\lambda^{q+1}+o(1), \quad \textrm{ as } k \to +\infty. 
\end{equation}
%
%
%Also, because of (\ref{eq:limitvkCaseA})-(ii)
%we get
%\begin{equation}\label{eq:case1concl3}
%	\lim_{k\to+\infty}\int_\Omega|v_k-v_\lambda|^{1-\gamma}=0.
%\end{equation}
Therefore, choosing $w=w_\lambda \in H_\lambda$ in \eqref{eq:ineq2.21haitao}, we obtain
\begin{equation}
	\begin{aligned}
		\|w_k - w_{\lambda}\|_{S^{1}_{0}(\Omega)}^2 &= -\int_{\Omega} \langle \nabla_{\mathbb{G}}w_k, \nabla_{\mathbb{G}}(w_{\lambda}-w_{k})\rangle_{\mathfrak{g}_1} + \int_{\Omega} \langle \nabla_{\mathbb{G}}w_{\lambda}, \nabla_{\mathbb{G}}(w_{\lambda}-w_{k})\rangle_{\mathfrak{g}_1}\\
		&\leq \dfrac{1}{k}\|w_{\lambda}-w_{k}\|_{S^{1}_{0}(\Omega)} + \int_{\Omega}w_{k}^{2\star_{Q}-1}(w_k - w_{\lambda}) + \lambda \int_{\Omega}w_{k}^{q}(w_k - w_{\lambda})\\
		&\qquad + \int_{\Omega} \langle \nabla_{\mathbb{G}}w_{\lambda}, \nabla_{\mathbb{G}}(w_{\lambda}-w_{k})\rangle_{\mathfrak{g}_1}\\
		&=\int_{\Omega}w_{k}^{2^{\star}_{Q}-1}(w_{\lambda}-w_k) + \lambda \int_{\Omega}w_{k}^{q+1} - \lambda \int_{\Omega}w_{k}^{q}w_{\lambda} +o(1)\\
		&=\|w_k-w_{\lambda}\|^{2^{\star}_{Q}}_{L^{2^{\star}_{Q}-1}(\Omega)} + \|w_{\lambda}\|^{2^{\star}_{Q}}_{L^{2^{\star}_{Q}-1}(\Omega)}-\int_{\Omega}w_{k}^{2^{\star}_{Q}-1}w_{\lambda} \\
		&\qquad + \lambda \int_{\Omega}w_{k}^{q+1} - \lambda \int_{\Omega}w_{k}^{q}w_{\lambda} +o(1)\\
		&=\|w_k-w_{\lambda}\|^{2^{\star}_{Q}}_{L^{2^{\star}_{Q}-1}(\Omega)} + o(1) \quad \textrm{ as } k \to +\infty.
	\end{aligned}
\end{equation}


%\[
%\begin{split}
%	\int_\Omega|\nabla_\G(v_k-v_\lambda)|^2+\lambda\int_\Omega v_k^{-\gamma}v_\lambda
%	&\leq \lambda\int_\Omega v_k^{1-\gamma}+\int_\Omega v_k^{2^\star_Q-1}(v_k-v_\lambda)+o(1)\\
%	&\leq \lambda\int_\Omega v_\lambda^{1-\gamma}+\|v_k-v_\lambda\|_{L^{2^\star_Q}(\Omega)}^{2^\star_Q}+o(1).	
%\end{split}
%\]
%Observing that $0\leq v_k^{-\gamma}v_\lambda\leq u_\lambda^{-\gamma}v_\lambda\in L^1(\Omega)$
%we can use again the Dominated Convergence Theorem,
%\begin{equation}
%	\int_\Omega v_k^{-\gamma}v_\lambda\to\int_\Omega v_\lambda^{1-\gamma}\ \ \m{as}\ k\to+\infty.
%\end{equation}
\noindent To proceed further, we choose $w=2w_k\in H_\lambda$,
yielding
\begin{equation}\label{eq:usando_2wk}
\|w_k\|_{S^{1}_{0}(\Omega)}^2-\|w_k\|_{L^{2^\star_Q}(\Omega)}^{2^\star_Q}-\lambda\int_\Omega w_k^{q+1}\geq-\frac{1}{k}\|w_k\|_{S^{1}_{0}(\Omega)}^2=o(1).
\end{equation}
Since $w_{\lambda}$ is actually a solution of \eqref{EqProblem}, we get
\begin{equation}\label{eq:wlambda_is_sol}
\|w_\lambda\|_{S^{1}_{0}(\Omega)}^2-\|w_\lambda\|_{L^{2^\star_Q}(\Omega)}^{2^\star_Q}-\lambda\int_\Omega w_k^{1-\gamma}\geq o(1)\ \ \m{as}\ k\to+\infty.
\end{equation}
Now, combining \eqref{eq:usando_2wk} with \eqref{eq:wlambda_is_sol} yields
 \begin{equation}\label{eq:vlambda1}
 	\|w_k-w_\lambda\|_{S^{1}_{0}(\Omega)}^2\geq \|w_k-w_\lambda\|_{L^{2^\star_Q}(\Omega)}^{2^\star_Q}+o(1)\ \ \m{as}\ k\to+\infty.
 \end{equation}
 Assuming without loss of generality that $I_\lambda(u_\lambda)\leq I_\lambda(w_\lambda)$,
 from \eqref{eq:EkelandCaseA} and \eqref{eq:case1concl1}-\eqref{eq:case1concl2}
 we obtain
 \[
 \begin{split}
 	I_\lambda(w_k-w_\lambda)&=I_\lambda(w_k)-I_\lambda(w_\lambda)+o(1)\\
 	&\leq I_\lambda(u_\lambda)-I_\lambda(w_\lambda)+\frac{1}{k^2}+o(1)\\
 	&=o(1)\ \ \m{as}\ k\to+\infty,
 \end{split}
 \]
 which, together with \eqref{eq:case1concl1}, gives
 \begin{equation}\label{eq:vlambda2}
 	\frac{1}{2}\|w_k-w_\lambda\|_{S^{1}_{0}(\Omega)}^2-\frac{1}{2^\star_Q}\|w_k-w_\lambda\|^{2^\star_Q}_{L^{2^\star_Q}(\Omega)}=
 	I_\lambda(w_k-w_\lambda)+\frac{\lambda}{q+1}\int_\Omega|w_k-w_\lambda|^{q+1}\leq o(1).
 \end{equation}
 From \eqref{eq:vlambda1} and \eqref{eq:vlambda2}
 we finally conclude
 \[
 \lim_{k\to+\infty}\|w_k-w_\lambda\|_{L^{2^\star_Q}(\Omega)}^{2^\star_Q}=\lim_{k\to+\infty}\|w_k-w_\lambda\|_{S^{1}_{0}(\Omega)}^2=0,
 \]
 proving \eqref{eq:stronconv}. This closes the proof.
\end{proof}
 
 
Combining Lemma \ref{lem:casoa1} with Lemma \ref{lem:casoa2} yields the existence of a second 
solution $w_\lambda$ to \eqref{EqProblem} for every $\lambda \in (0,\Lambda)$ in the case that
for every $r\in (0,r_0)$,
\[
\inf_{\|u-u_\lambda\|_{S_{0}^{1}(\Omega)}=r}I_\lambda(u)=I_\lambda(u_\lambda).
\]
 

\subsection{Second case} \label{subsec:CaseII}

%We are going to prove that if there exists $r_1\in(0,r_0)$ such that
%\begin{equation}\label{eq:case2}
%\inf_{\|u-u_\lambda\|=r_1}I_\lambda(u)>I_\lambda(u_\lambda)
%\end{equation}
%for some $r_1\in (0,r_0)$, then 
%there exists a (second) solution $v_\lambda$ of the problem \eqref{EqProblem}$_\lambda$
%such that $0<u_\lambda<v_\lambda$.\newline

%\color{red}
%As before, in the case $p\in (1,2^*_Q-1)$, the functional $I_\lambda$
%respects the PS condition. Moreover, we show (BY FOLLOWING ADV)
%that there exists $z_\lambda\in S^1_0(\Omega)$ such that
%\[
%I_\lambda(z_\lambda)<I_\lambda(u_\lambda).
%\]
%\color{black}
%This ensures a Mountain-Pass geometry of the problem,
%and therefore the existence of a critical point $w_\lambda\neq u_\lambda$.\newline
%

%\color{red} We can focus again in the case $p=2^*_Q-1$,
%where we recover another critical point
%even without the compactness hypothesis.\color{black}  

As before, we consider the space $H_\lambda$ introduced in \eqref{def:hlambda} and
$C([0,1],H_\lambda)$ the space of continuous curves endowed with the following distance,
\begin{equation}\label{maxdistance}
d(\eta,\eta')=\max_{t\in[0,1]}\|\eta(t)-\eta'(t)\|_{S^1_0(\Omega)}.
\end{equation}
Moreover, we consider the following subspace
\begin{equation}\label{def:gammal}
\Gamma_\lambda:=\left\{\eta\in C([0,1],H_\lambda):\begin{array}{l}
	%\diamond&
	 \eta(0)=u_\lambda,\\[0.1cm] 
	%\diamond& 
	\|\eta(1)-u_\lambda\|_{S^{1}_{0}(\Omega)}>r_1,\\[0.1cm]
	%\diamond& 
	I_\lambda(\eta(1))<I_\lambda(u_\lambda)
\end{array}
\right\}.
\end{equation}

\noindent We first show that $\Gamma_\lambda\neq \varnothing$
and then we provide an estimate of the minimax level
\[
\ell_0:=\inf_{\eta\in\Gamma_\lambda}\max_{t\in[0,1]}I_\lambda(\eta(t)).
\]

\noindent Once this is done, we will apply once again the Ekeland's variational principle.

\noindent We consider the family of functions $U_\varepsilon$ defined
at \eqref{eq:def_u_varepsilon}.
For any $a\in \G$, we take
\[
U_{\varepsilon,a}(g):=U_\varepsilon(a^{-1}\diamond g)
= {\varphi(a^{-1}\diamond g) T_{\varepsilon}(a^{-1}\diamond g)}.
\]
We recall that $\{T_\varepsilon\}$ is a family
of functions defined in \eqref{eq:Minimi_riscalati} such that  $T_1=T$
is a minimizer of the Sobolev Inequality \eqref{eq:Sobolev_Ineq}.
%where the dot stands for the multiplication in $\G$.


\begin{lemma}\label{lem:claim2.7haitao}
	There exists $\varepsilon_0>0$, $a\in\Omega$ and $R_0\geq1$ such that
\begin{align}
	 & I_\lambda(u_\lambda + R U_{\varepsilon,a}) < I_\lambda(u_\lambda) 
	&& \forall \varepsilon \in (0, \varepsilon_0),\ \forall R \geq R_0, \\
& I_\lambda(u_\lambda + t R_0 U_{\varepsilon,a}) < I_\lambda(u_\lambda) + \frac{1}{Q} S_\G^{Q/2} 
	&& \forall t \in [0, 1],\ \forall \varepsilon \in (0, \varepsilon_0),\label{eq:lemma2.7Haitao}
\end{align}
where $S_\G$ is the best Sobolev constant defined as in  \eqref{eq:def_best_constant}.
\end{lemma}
\begin{proof}
%We set the notation $p:=2^\star_Q-1=\frac{Q+2}{Q-2}$.
Following the computations already developed in \cite{BiGaVe},
\begin{align*}
    I_\lambda(u+tRU_{\varepsilon,a})&=
    \frac{1}{2}\int_{\Omega}|\nabla_{\G}u|^2-\frac{1}{2^\star_Q}\int_{\Omega} |u|^{2^\star_Q}-\frac{\lambda}{q+1}\int_{\Omega} |u|^{q+1}\\
    &\ \ \ \ +tR \left(\int_{\Omega}\langle \nabla_{\G}u,\nabla_{\G}U_{\varepsilon,a}\rangle_{\mathfrak{g}_{1}}-\int_{\Omega} u^{2^\star_Q -1}U_{\varepsilon,a}-\lambda\int u^qU_{\varepsilon,a}\right)\tag{C}\label{eq:C}\\
    &\ \ \ \ -\frac{1}{2^\star_Q}\left(\int|u+tRU_{\varepsilon,a}|^{2^\star_Q}-\int_{\Omega} |u|^{2^\star_Q}\right)+tR\, \int_{\Omega} u^{2^\star_Q -1}U_{\varepsilon,a}\tag{A}\label{eq:A1}\\
     &\ \ \ \ -\frac{\lambda}{q+1}\left(\int_{\Omega}|u+tRU_{\varepsilon,a}|^{q+1}-\int_{\Omega} |u|^{q+1}\right)+ \lambda tR\, \int_{\Omega} u^qU_{\varepsilon,a}\tag{A2}\label{eq:A2}\\
    &\ \ \ \ +\frac{t^2R^2}{2} \int_{\Omega}|\nabla_{\G}U_{\varepsilon,a}|^2.\tag{B}\label{eq:B}
    %&\ \ \ \ -\frac{\lambda}{1-\gamma}\left(\int(u+tRU_{\varepsilon,a})^{1-\gamma}-\int u^{1-\gamma}\right)+\lambda\, tR\, \int u^{-\gamma}U_{\varepsilon,a}\tag{D}\label{eq:D}
\end{align*}
Taking the limit for $\varepsilon\to0^+$, we 
start by estimating~\eqref{eq:A2}.
We observe that 
\[
tRu^{2^\star_Q -1}U_{\varepsilon,a}=\left.\left((u+s)^{2^\star_Q}\right)'\right|_{s=tRU_{\varepsilon,a}}.
\]
Therefore, by  strict convexity the term \eqref{eq:A2} must be negative.

For the other terms, we take analogous estimates to the ones developed
in \cite{BiGaVe}. In particular, \eqref{eq:C} is null by definition of
weak solution with $U_{\varepsilon,a}$ as test function,
while the other terms have been estimated in \cite[Lemma 4.3]{BiGaVe}.
We resume,
\begin{align*}
    \eqref{eq:C}&\ =0,\\
\eqref{eq:A1}&\ =-\frac{A\, t^{2^\star_Q}R^{2^\star_Q}}{2^\star_Q}-t^{2^\star_Q -1}R^{2^\star_Q -1}\, K\, \varepsilon^{\frac{Q-2}{2}}+o\left(\varepsilon^{\frac{Q-2}{2}}\right),\\
%\eqref{eq:A2}&\ =-\frac{A_2\, t^{q+1}R^{q+1}}{q+1}-t^qR^q\, K\, \varepsilon^{\frac{Q-2}{2}}+o\left(\varepsilon^{\frac{Q-2}{2}}\right)\\
    \eqref{eq:B}&\ =\frac{B\, t^2R^2}{2}+o\left(\varepsilon^{\frac{Q-2}{2}}\right),
%    \eqref{eq:D}&\ =o\left(\varepsilon^{\frac{Q-2}{2}}\right)
\end{align*}
for suitable positive constants $A,B,K$.
In particular,
\begin{align*}
	A&=\|T\|_{L^{2^\star_Q}(\G)},\\
	B&=\||\nabla_\G T|\|_{L^2(\G)}^2,
\end{align*}
and both can be estimated with the Sobolev constant
as done in Lemma \ref{lem:Stime_Talentiane_modificate}.
Moreover, $K$ is a constant depending on $a$ with the property that
\[
\int_\Omega uU_{\varepsilon,a}^{2^\star_Q -1}=K\varepsilon^{(Q-2)\slash 2}+o(\varepsilon^{(Q-2)\slash2}).
\]

To conclude the proof of the lemma 
we follow the approach of \cite{Tarantello}.
We slightly change the notation by posing
 $s:=tR$ and $S:=\left(\frac{B}{A}\right)^{\frac{1}{2^\star_Q -2}}$,
and, we introduce
\[
f_\varepsilon(s):=\frac{B\, s^2}{2}-\frac{A \, s^{2^\star_Q}}{2^\star_Q}-s^{2^\star_Q -1}\, K\, \varepsilon^n.
\]
As \eqref{eq:A2} is negative, we have
\begin{equation}\label{disIf}
	I_\lambda(u+tRU_{\varepsilon,a})<f_\varepsilon(tR)+o\left(\varepsilon^{\frac{Q-2}{2}}\right).
\end{equation}

With the same calculations developed in \cite{BiGaVe}
we get
%
%
%We denote by $s_\varepsilon>0$ the point where $f_\varepsilon$ achieves its max. Observe that
%\begin{equation}\label{eqboh}
%	f'_\varepsilon(s)=Bs-As^p-ps^{p-1}K\varepsilon^n
%\end{equation}
%and also
%\begin{equation}\label{bsas}
%BS-AS^p=0.
%\end{equation}
%The last equations \eqref{eqboh} and \eqref{bsas} imply
%\[
%S>s_\varepsilon>0\ \m{ and } \ s_\varepsilon\to S\m{ as }\varepsilon\to0^+.
%\]
%Let's write
%\[
%s_\varepsilon=S(1-t_\varepsilon).
%\]
%We can develop for $t_\varepsilon$,
%\[
%t_\varepsilon\, (p-1)\left(\frac{B^p}{A}\right)^{\frac{1}{p-1}}=p\, \frac{B}{A}\, K\, \varepsilon^n+o(\varepsilon^n), \quad \textrm{ as } \varepsilon \to 0^+.
%\]
%and we obtain, as $\varepsilon \to 0^+$,
%\begin{align*}
%    I_\lambda(u+tRU_{\varepsilon,a})
%    &<I_\lambda(u)+\frac{Bs_\varepsilon^2}{2}-\frac{As_\varepsilon^{p+1}}{p+1}-s_\varepsilon^{p}K\varepsilon^n+o(\varepsilon^n)\\
%    &=I_\lambda(u)+\frac{BS^2}{2}-\frac{AS^{p+1}}{p+1}-BSt_\varepsilon+AS^pt_\varepsilon-S^{p}K\varepsilon^n+o(\varepsilon^n)\\
%    &\ (\m{and therefore by Eq.~\eqref{bsas}})\\
%    &=I_\lambda(u)+\frac{BS^2}{2}-\frac{AS^{p+1}}{p+1}-S^pK\varepsilon^n+o(\varepsilon^n)\\
%    &=I_\lambda(u)+\left(\frac{1}{2}-\frac{1}{p+1}
%    \right)\, \frac{B^{\frac{p+1}{p-1}}}
%    {A^{\frac{2}{p-1}}}
%    -S^pK\varepsilon^n+o(\varepsilon^n)
%\end{align*}
\[
I_\lambda(u+tRU_{\varepsilon,a})<I_\lambda(u)+\left(\frac{1}{2}-\frac{1}{2^\star_Q}\right)\cdot \frac{B^{\frac{2^\star_Q}{2^\star_Q -2}}}{A_1^{\frac{2}{2^\star_Q -2}}}-S^{2^\star_Q -1}K\varepsilon^{\frac{Q-2}{2}}
+o\left(\varepsilon^{\frac{Q-2}{2}}\right),
\]
and this allows to conclude because $\frac{1}{2}-\frac{1}{2^\star_Q}=\frac{1}{Q}$
and $\frac{B}{A_1^\frac{2}{2^\star_Q}}=S_{\G}$. This closes the proof.
\end{proof}

We notice that by Lemma \ref{lem:claim2.7haitao}, $\Gamma_\lambda\neq \varnothing$ and moreover $\Gamma_\lambda$ is a complete metric space endowed
with the distance~\eqref{maxdistance}. This follows simply by the fact that
\begin{equation}\label{eq:eta}
	\widetilde \eta(t) := u_\lambda+tR_0U_{\varepsilon,a}\in \Gamma_\lambda\quad\text{for all $\varepsilon\in (0,\varepsilon_0)$},
\end{equation}
\noindent eventually enlarging $R_0$.
We now proceed by showing that problem \eqref{EqProblem} do actually admit a second solution.

\begin{teo}\label{thm:casob}
	For $0<\lambda<\Lambda$, if there exists $r\in (0,r_0)$ such that 
	\[
	\inf_{\|u-u_\lambda\|_{S_{0}^{1}(\Omega)}=r}I_\lambda(u)>I_\lambda(u_\lambda),
	\]
	then there exists a solution $v_\lambda$ to \eqref{EqProblem}
	such that $v_\lambda \not\equiv u_\lambda$.\newline
\end{teo}
\begin{proof}
We start by considering the Generalized Directional Derivative
introduced in \cite{BadTar}, 
\[
I^0_\lambda(w;v):=\limsup_{\|h\|\to0^+,\, \rho\to0^+}\frac{I_\lambda(w+h+\rho v)-I_\lambda(w+h)}{\rho}.
\]
Some basic properties of this object are proved in \cite[Appendix]{BiGaVe}.
In the case of object,
\[
I_\lambda^0(w;v)=\int_\Omega\langle \nabla_\G w,\nabla_\G v\rangle_{\mathfrak{g}_1}-\lambda\int_\Omega w^{q}v-\int_\Omega w^{2^\star_Q-1}v.
\]
We recall the functional
\[
\Phi(\eta):= \max_{t\in [0,1]}I_\lambda(\eta(t)).
\]
and its minimax level
\[
\ell_0:=\inf_{\eta\in\Gamma_\lambda}\Phi(\eta).
\]
By Ekeland's Variational Principle, there exists
a sequence $\{\eta_k\}_k\in \Gamma_\lambda$
such that
$$\Phi(\eta_k)\leq \ell_0+\frac{1}{k} \quad \textrm{ and } \quad \Phi(\eta_k)\leq \Phi(\eta)+\frac{1}{k}d(\eta_k,\eta).$$
We now recall the result of \cite[Lemma A.2]{BiGaVe}.
For every $k\in \mathbb{N}$, we define 
\[
\Lambda_k:=\left\{t\in(0,1):\ I_\lambda(\eta_k(t))=\max_{s \in [0,1]}I_\lambda(\eta_k(s))\right\}.
\]
Then, for every $k$ there exists $t_k\in \Lambda_k$ such that 
for $v_k:=\eta_k(t_k)$,
\begin{equation}\label{lemmaimp}
	I^0_\lambda(v_k;w-v_k)\geq-\frac{\max(1,\|w-v_k\|)}{k}\quad\textrm{ for every } w\in H_\lambda.
\end{equation}
Resuming what we just presented,
\begin{itemize}
	%\item $v_k:=\eta_k(t_k)\in H_\lambda$
	\item $I_\lambda(v_k)\to\ell_0$ as $k\to+\infty$,
	\item there exists $C>0$ such that $\forall w\in H_\lambda$,
	\begin{equation}\label{eq:badtarfin}
		\int_\Omega\langle \nabla_\G v_k,\nabla_\G(w-v_k)\rangle_{\mathfrak{g}_{1}} -\lambda\int_\Omega v_k^{q}(w-v_k)-\int_\Omega v_k^{2^\star_Q-1}(w-v_k)\geq
		-\frac{C}{k}(1+\|w\|_{S^{1}_{0}(\Omega)}).
	\end{equation}
\end{itemize}
By choosing $w=2v_k$,
we get
\begin{equation}\label{eq:ultima}
\| v_k\|^2_{{S_0^1(\Omega)}}-\|v_k\|^{2^\star_Q}_{L^{2^\star_Q}(\Omega)}-\lambda \|v_k\|^{q+1}_{L^{q+1}(\Omega)}\geq-\frac{C}{k}\max(1,\|v_k\|_{{S_0^1(\Omega)}}).
\end{equation}

\noindent By \eqref{eq:ultima}, and recalling \eqref{eq:Def_Ilambda}, it follows that
\begin{equation*}
	\begin{aligned}
	\dfrac{1}{2}&\| v_k\|^2_{{S_0^1(\Omega)}} - \dfrac{\lambda}{q+1}\|v_k\|^{q+1}_{L^{q+1}(\Omega)} - \dfrac{1}{2^\star_Q}\|v_k\|^{2^\star_Q}_{L^{2^\star_Q}(\Omega)} \\
	&-\dfrac{1}{2^{\star}_{Q}}\left(\| v_k\|^2_{{S_0^1(\Omega)}}-\|v_k\|^{2^\star_Q}_{L^{2^\star_Q}(\Omega)}-\lambda \|v_k\|^{q+1}_{L^{q+1}(\Omega)} +\frac{C}{k}\max(1,\|v_k\|_{{S_0^1(\Omega)}}) \right) \leq I_{\lambda}(v_k).
	\end{aligned}
\end{equation*}


\noindent As said above, $I_\lambda(v_k)\to \ell_0$ for $k\to+\infty$,
and therefore we obtain
%\[
%I_\lambda(v_k)-\frac{1}{2^\star_Q}\cdot \left(\m{Equation }\eqref{eq:ultima}\right)
%\]

\begin{equation}\label{eq:gamma0}
	\ell_0+o(1)\geq \left(\frac{1}{2}-\frac{1}{2^\star_Q}\right)\|v_k\|_{{S_0^1(\Omega)}}^2-\lambda\left(\frac{1}{q+1}-\frac{1}{2^\star_Q}\right)\|v_k\|_{L^{q+1}(\Omega)}^{q+1}.
\end{equation}

% the equation above, we get
%\begin{equation}\label{eq:last}
%	\|v_k\|_{S^{1}_{0}(\Omega)}^2-\lambda\int_\Omega v_k^{1-\gamma}-\int_\Omega v_k^{2^\star_Q}\geq
%	-\frac{C}{k}(1+2\|v_k\|_{S^{1}_{0}(\Omega)}).
%\end{equation}
%Using H\"older's and Sobolev's inequality,
%and the fact that $I_\lambda(v_k)\to \ell_0$, we get
%\begin{equation}\label{eq:gamma0}
%	\begin{split}
%		\ell_0+o(1)
%		&= \frac{1}{2}\|v_k\|_{S^{1}_{0}(\Omega)}^2-\frac{\lambda}{1-\gamma}\int_\Omega v_k^{1-\gamma}
%		-\frac{1}{2^\star_Q}\int_\Omega v_k^{2^\star_Q}\\
%		&\geq\left(\frac{1}{2}-\frac{1}{2^\star_Q}\right)\|v_k\|_{S^{1}_{0}(\Omega)}^2-\lambda\left(\frac{1}{1-\gamma}-\frac{1}{2^\star_Q}\right)\int_\Omega v_k^{1-\gamma}
%		-\frac{C}{2^\star_Q\, k}(1+2\|w\|)\\
%		&\geq \left(\frac{1}{2}-\frac{1}{2^\star_Q}\right)\|v_k\|_{S^{1}_{0}(\Omega)}^2-C(\|v_k\|_{S^{1}_{0}(\Omega)}^{1-\gamma}-2\|v_k\|_{S^{1}_{0}(\Omega)}-1).
%	\end{split}
%\end{equation}
%We recall that the constant $C$ depends on $Q$ and $|\Omega|$.

Suppose $v_k$ is unbouded in $S^1_0(\Omega)$. As
$\frac{1}{2}-\frac{1}{2^\star_Q}>0$, 
then up to a subsequence we have
 $\|v_k\|_{S^{1}_{0}(\Omega)}\to+\infty$,
contradicting \eqref{eq:gamma0}.
Therefore $v_k$ is bounded
and \eqref{lemmaimp} with
$w=2v_k$ implies
\begin{equation}\label{eq:last}
	\|v_k\|_{S^{1}_{0}(\Omega)}^2-\lambda\int_\Omega v_k^{q+1}-\int_\Omega v_k^{2^\star_Q}\geq
	-\frac{C}{k}(1+2\|v_k\|_{S^{1}_{0}(\Omega)}).
\end{equation}

Now, we can proceed following almost verbatim \cite{BiGaVe}. Arguing as in 
the proofs of Lemmas \ref{lem:casoa1} and \ref{lem:casoa2},
we can prove both that
$v_k$ weakly converges (up to a subsequence) to a weak solution $v_\lambda$
of \eqref{EqProblem} and that
\begin{equation}\label{eq:convergence}
	\|v_k-v_\lambda\|_{S^{1}_{0}(\Omega)}^2-\|v_k-v_\lambda\|_{L^{2^\star_Q}(\Omega)}^{2^\star_Q}=o(1)\ \ \m{as}\ k\to+\infty.
\end{equation}

It only remains to show that $v_\lambda\not \equiv u_\lambda$, i.e. that we truly found a second solution of \eqref{EqProblem}.
To this aim, we notice first that for any $\eta\in \Gamma_\lambda$,
\[
\|\eta(0)-u_\lambda\|_{S^{1}_{0}(\Omega)}=0\ \ \m{and}\ \ \|\eta(1)-u_\lambda\|_{S^{1}_{0}(\Omega)}>r_1,
\]
therefore there exists $t_\eta\in [0,1]$
such that $\|\eta(t_\eta)-u_\lambda\|_{S^{1}_{0}(\Omega)}=r_1$.
Since we are dealing with (2),
we have
\begin{equation}\label{eq:Stima_l0_basso}
\ell_0=\inf_{\Gamma_\lambda}\Phi(\eta)\geq\inf_{\Gamma_\lambda}I_\lambda(\eta(t_\eta))
\geq\inf_{\|u-u_\lambda\|_{S^{1}_{0}(\Omega)}=r_1} I_\lambda(u)>I_\lambda(u_\lambda),
\end{equation}
\noindent where the last infimimum is taken among the $u\in H_\lambda$
such that $\|u-u_\lambda\|_{S^{1}_{0}(\Omega)}=r_1$. On the other hand
if we consider $\widetilde\eta$ defined in \eqref{eq:eta},
then by \eqref{eq:lemma2.7Haitao}
we get,
\begin{equation}\label{eq:Stima_l0_alto}
\ell_0\leq \Phi(\widetilde \eta)=\max_{t\in [0,1]}I_\lambda(\widetilde \eta(t))<
I_\lambda(u_\lambda)+\frac{1}{Q}S_\G^{Q\slash2}.
\end{equation}
Combining \eqref{eq:Stima_l0_basso} with \eqref{eq:Stima_l0_alto} we find
\begin{equation}\label{eq:sumup}
	I_\lambda(u_\lambda)<\ell_0<I_\lambda(u_\lambda)+\frac{1}{Q}S_\G^{Q\slash2}.
\end{equation}
We further stress that, since $v_k\to v_\lambda$ weakly in $S^1_0(\Omega)$, then
the equalities \eqref{eq:case1concl1}-\eqref{eq:case1concl2} still hold true.
Therefore, by \eqref{eq:sumup} and recalling that $I_\lambda(v_k)\to \ell_0$,
we get (for $k$ sufficiently large), that there exists a positive number $\delta_{0}$ such that 
\begin{equation}\label{eq:sumup2}
	\begin{split}
		\frac{1}{2}&\|v_k-v_\lambda\|_{S^{1}_{0}(\Omega)}^2-\frac{1}{2^\star_Q}\|v_k-v_\lambda\|_{L^{2^\star_Q}(\Omega)}^{2^\star_Q}\\
		&=\frac{1}{2}(\|v_k\|_{S^{1}_{0}(\Omega)}^2-\|v_\lambda\|_{S^{1}_{0}(\Omega)}^2)-\frac{1}{2^\star_Q}(\|v_k\|^{2^\star_Q}_{L^{2^\star_Q}(\Omega)})+o(1)\\
		&=I_\lambda(v_k)-I_\lambda(u_\lambda)+o(1)\\
		&=\ell_0-I_\lambda(u_\lambda)+o(1)\\
		&<\frac{1}{Q}S_\G^{Q\slash2}-\delta_0,
	\end{split}
\end{equation}
\noindent and the last term is positive.
From \eqref{eq:convergence}, \eqref{eq:sumup} and \eqref{eq:sumup2},
closely following \cite[Proposition 3.1]{Tarantello},
we get that $v_k\to v_\lambda$ strongly in $S^1_0(\Omega)$.
This, also considering both \eqref{eq:last} and \eqref{eq:gamma0}, gives
\[
I_\lambda(u_\lambda)<\gamma_0=\lim_{k\to+\infty}I_\lambda(v_k)=I_\lambda(v_\lambda),
\]
thus implying that $u_\lambda\not\equiv v_\lambda$. This closes the proof.
\end{proof}

%We resume all this in the following lemma.
%\begin{lemma}\label{lem:casob}
%	 For $0<\lambda<\Lambda$, if there exists $r\in (0,r_0)$ such that 
%	\[
%	\inf_{\|u-u_\lambda\|_{S_{0}^{1}(\Omega)}=r}I_\lambda(u)>I_\lambda(u_\lambda),
%	\]
%	then there exists a solution $v_\lambda$ to \eqref{EqProblem}$_{\lambda}$
%	such that $v_\lambda \not\equiv u_\lambda$.\newline
%\end{lemma}
\medskip


Gathering all the results established so far, we can finally provide the
\begin{proof}[Proof of Theorem \ref{thm:main}]
	Let $\Lambda$ be as in \eqref{eq:DefinitionLambda}, that is, 
    $$\Lambda := \sup \{ \lambda >0: \eqref{EqProblem} \textrm{ admits a weak solution}\}.$$
    From its very definition, this shows that \eqref{EqProblem}
    does not admit any weak solution for $\lambda > \Lambda$, which is assertion A).
    Moreover, combining Lemma \ref{lem:lambda0} with Lemma \ref{lem:LambdaInf}, we immediately have that $\Lambda\in (0,+\infty)$.
    Regarding the existence part, Theorem \ref{thm:Existence_First} proves that 
    there exists \emph{at least one weak solution} $u_\lambda$ of
    \eqref{EqProblem} (with $\lambda =\Lambda$), which is assertion B).
    As for assertion C),
    from Lemmas  \ref{lem:casoa1}, \ref{lem:casoa2}  and Theorem \ref{thm:casob},
    we got that a second
    solution always exists when $0<\lambda<\Lambda$.
    % it is enough to take into account the
    %computations performed in Sections \ref{subsec:CaseI}-\ref{subsec:CaseII}. This closes the proof.
\end{proof}
\medskip


\begin{thebibliography}{99}
	
	
	\bibitem{AbDiVa}
	B. Abdellaoui, A. Dieb, E. Valdinoci,
	{\em A nonlocal concave-convex problem with nonlocal mixed boundary data},
	Commun. Pure Appl. Anal. {\bf 17}(3), (2018), 1103--1120.
	
	\bibitem{AbTr}
	F. Abedin, G. Tralli,
	{\em Boundary regularity for subelliptic equations in the Heisenberg group},
	preprint. \url{https://arxiv.org/abs/2506.05151}
	
%	\bibitem{Afeltra}
%	C. Afeltra, 
%	{\em Singular periodic solutions to a critical equation in the Heisenberg group},
%	Pacific J. Math. {\bf 305}(2), (2020), 385--406.
	
	\bibitem{Alama}
	S. Alama,
	{\em Semilinear elliptic equations with sublinear indefinite nonlinearities},
	Adv. Differential Equations {\bf 4}, (1999), 813--842.
	
%	\bibitem{Alves}
%	C.O. Alves, S. Gandal, A. Loiudice, J. Tyagi, 
%	{\em A Br\'{e}zis-Nirenberg type problem for a class of degenerate elliptic problems involving the Grushin operator},	
%	J. Geom. Anal. {\bf 34}(2), (2024), Paper No. 52, 41 pp.
%	
\bibitem{ABC}
   A. Ambrosetti, H. Br\'{e}zis, G. Cerami, 
   {\em Combined effects of concave and convex nonlinearities in some elliptic problems},
   J. Funct. Anal. {\bf 122}(2), (1994), 519--543.


%	\bibitem{ArNgRa}
%	R. Arora, P.-T. Nguyen, V.D. Radulescu,
%	{\em A large class of nonlocal elliptic equations with singular nonlinearities},
%	preprint. \url{https://arxiv.org/abs/2211.06634}
	
	
	%	\bibitem{ArRa}
	%	R. Arora, V.D. Radulescu,
	%	\emph{Combined effects in mixed local-nonlocal stationary problems}, 
	%	Proc. Roy. Soc. Edinburgh Sect. A, (2023), pp. 1--47. \url{https://doi.org/10.1017/prm.2023.80}
	
%	\bibitem{AubEke}
%	J.-P. Aubin, I. Ekeland, 
%	{\em Applied nonlinear analysis}.
%	Pure Appl. Math. (N.Y.) Wiley-Intersci. Publ.
%	John Wiley \& Sons, Inc., New York, 1984. xi+518 pp.
	
	\bibitem{BadTar}
	M. Badiale, G. Tarantello, 
	{\em Existence and multiplicity results for elliptic problems with critical growth and discontinuous nonlinearities}, 
	Nonlinear Anal. Theory Methods Appl. {\bf 29}, (1997), 639--677. 
	
	\bibitem{BaCiCu}
	A. Baldi, G. Citti, G. Cupini, 
	{\em Schauder estimates at the boundary for sublaplacians in Carnot groups}, 
	Calc. Var. Partial Differential Equations {\bf 58} (2019), art. 204.
	
	\bibitem{BaGaMu}
	A. Banerjee, N. Garofalo, I.H. Munive, 
	{\em Higher order boundary Schauder estimates in Carnot groups}, 
	Math. Ann. {\bf 390} (2024), 6013--6047.
	
	\bibitem{BESS}
	B. Barrios, E. Colorado, R. Servadei, F. Soria, 
	{\em A critical fractional equation with concave-convex
		power nonlinearities}, 
	Ann. Inst. H.Poincar\'{e} Anal. Non Lin\'{e}aire {\bf 32}(4), (2015), 875--900.
	
	%	\bibitem{BDBMP}
	%	B. Barrios, I. De Bonis, M. Medina, I. Peral, 
	%	{\em Semilinear problems for the fractional laplacian with a singular nonlinearity}, 
	%	Open Math. {\bf 13}, (2015), 390--407.
	
	\bibitem{BiGaVe}
	S. Biagi, M. Galeotti, E. Vecchi,
	{\em Critical singular problems in Carnot groups}, preprint. \url{https://arxiv.org/abs/2506.07521}
	
	\bibitem{BPV}
	S. Biagi, A. Pinamonti, E. Vecchi, 
	{\em Sublinear Equations Driven by Hörmander Operators}, J. Geom. Anal. {\bf 32}, (2022), art. 121.
	
	\bibitem{BV} 
	S. Biagi, E. Vecchi,
	{\em On the existence of a second positive solution to mixed local-nonlocal concave-convex critical problems}, Nonlinear Anal. {\bf 256}, (2025), 113795. 
	
%	\bibitem{BiCa}
%	I. Birindelli, I. Capuzzo Dolcetta, A. Cutrì, 
%	{\em Liouville theorems for semilinear equations on
%		the Heisenberg group}, 
%	Ann. Inst. H. Poincar\'{e} C Anal. Non Lin\'{e}aire {\bf 14}, (1997), 295--308.
	
%	\bibitem{BiLa}
%	I. Birindelli, E. Lanconelli, 
%	{\em A negative answer to a one-dimensional symmetry problem in the Heisenberg group},
%	Calc. Var. Partial Differential Equations {\bf 18}(4), (2003), 357--372.
	
%	\bibitem{BiPr}
%	I. Birindelli, J.V. Prajapat, 
%	{\em Nonlinear Liouville theorems in the Heisenberg group via the moving plane method},	
%	Comm. Partial Differential Equations {\bf 24}(9-10), (1999), 1875--1890.
	
%	\bibitem{BiPr2}	
%	I. Birindelli, J.V. Prajapat,
%	{\em One dimensional symmetry in the Heisenberg group},
%	Ann. Scuola Normale Superiore di Pisa Cl. Sci. (5) {\bf 30}(4), (2001), 269--284.	
	
	%	\bibitem{BO}
	%	L. Boccardo, L. Orsina, 
	%	{\em Semilinear elliptic equations with singular nonlinearities}, Calc. Var. Partial Differ. Equ. {\bf 37}, (2010), 363--380.
	
	\bibitem{BLU}
	A. Bonfiglioli, E. Lanconelli, F. Uguzzoni, 
	{\em Stratified Lie Groups and Potential Theory for their
		Sub-Laplacians}. Springer Monographs in Mathematics, vol. 26. Springer, New York, NY (2007).
	
	\bibitem{BoUg}
	A. Bonfiglioli, F. Uguzzoni,
	{\em Nonlinear Liouville theorems for some critical
		problems on H-type groups},
	J. Funct. Anal. {\bf 207}(1), (2004), 161--215.
	
%	\bibitem{Bony}
%	J.M. Bony,
%	{\em Principe du maximum, in\'{e}galit\'{e} de Harnack et unicit\'{e} du probl\`{e}me de Cauchy pour les op\'{e}rateurs elliptiques d\'{e}g\'{e}n\'{e}r\'{e}s}, 
%	Ann. Inst. Fourier Grenobles {\bf 19}(1), (1969), 277--304.
	
%	\bibitem{BoFiPu}
%	S. Bordoni, R. Filippucci, P. Pucci, 
%	{\em Existence problems on Heisenberg groups involving Hardy and
%		critical terms}, 
%	J. Geom. Anal. {\bf 30}(2), (2020), 1887--1917.
	
	
	
%	\bibitem{BrRiSe}
%	L. Brandolini, M. Rigoli, A.G. Setti, 
%	{\em Positive solutions of Yamabe-type equations on the Heisenberg group},
%	Duke Math. J. {\bf 91}, (1998), 241--296.
%	
	%\bibitem{Brezis}
	%H. Brezis,
	%\emph{Functional analysis, Sobolev spaces and partial differential equations}, 
	% U\-ni\-ver\-si\-text, Springer, New York, 2011. 
	
	\bibitem{BrezisKamin}
	H. Br\'{e}zis, S. Kamin, 
	{\em Sublinear elliptic equations in $\mathbb{R}^{n}$},
	Manuscripta Math. {\bf 74}(1), (1992), 87--106.
	
%	\bibitem{BrezisLieb}
%	H. Br\'{e}zis, E. Lieb, 
%	{\em A relations between pointwise convergence of functions and convergence of integrals}, 
%	Proc. Amer. Math. Soc. {\bf 88}, (1983), 486--490.
	
	\bibitem{BN}
	H. Br\'{e}zis, L. Nirenberg, 
	{\em Positive solutions of nonlinear elliptic equation involving the critical Sobolev exponent}, 
	Comm. Pure Appl. Math. {\bf 36}, (1983), 437--477. 
	
	
%	\bibitem{BrezisNirenberg}
%	H. Br\'{e}zis, L. Nirenberg, 
%	{\em A minimization problem with critical exponent and nonzero data}, 
%	in Symmetry in Nature (a volume in honor of L. Radicati), Scuola Normale Superiore Pisa, (1989), Volume I, 129--140.
	
		\bibitem{BNH1C1}
		H. Br\'{e}zis, L. Nirenberg, 
		{\em $H^1$ versus $C^1$ local minimizers}, 
		C. R. Acad. Sci. Paris {\bf 317}, (1993), 465--472. 
	
		
%	\bibitem{CLMR}
%	G. Catino, Y. Li, D.D. Monticelli, A. Roncoroni,
%	{\em A Liouville theorem in the Heisenberg group},
%	to appear J. Eur. Math. Soc.
	
	\bibitem{CCP}
	F. Charro, E. Colorado, I. Peral, 
	{\em Multiplicity of solutions to uniformly elliptic fully nonlinear equations with concave-convex right-hand side}, 
	J. Differential Equations {\bf 246}(11), (2009), 4221--4248.
	
%	\bibitem{ChMaYa}
%	J.-H. Cheng, A. Malchiodi, P. Yang,
%	{\em On the Sobolev quotient
%	of three-dimensional CR manifolds},
%	Rev. Mat. Iberoam. {\bf 39}(6), (2023), 2017--2066.
	
	\bibitem{Citti}
	G. Citti,
	{\em Semilinear Dirichlet problem involving critical exponent for the Kohn Laplacian}, 
	Ann. Mat. Pura. Appl. {\bf 169}, (1995), 375--392.
	
	\bibitem{CitUg}
	G. Citti, F. Uguzzoni, 
	{\em Critical semilinear equations on the Heisenberg group: the effect of the
		topology of the domain}, Nonlinear Anal. {\bf 46}, (2001), 399--417.
	
	%\bibitem{CoifWeiss}
	%R.R. Coifman, G. Weiss, 
	%{\em Analyse Harmonique Non-commutative sur Certains Espaces Homogènes}. Lecture Notes in Math No. 242. Springer, Berlin (1971).
	
	\bibitem{CoPe}
	E. Colorado, I. Peral, 
	{\em Semilinear elliptic problems with mixed Dirichlet-Neumann boundary conditions}, 
	J. Funct. Anal. {\bf 199}(2), (2003), 468--507.

	
	%\bibitem{Cygan}
	%J. Cygan, 
	%{\em Subadditivity of homogeneous norms on certain nilpotent Lie groups}, 
	%Proc. Amer. Math. Soc. {\bf 83}(1), (1981), 69--70.
	
	%	\bibitem{DiBenedetto}
	%	E. DiBenedetto,
	%	{\em Degenerate parabolic equations},
	%	Universitext, Springer-Verlag, New York, 1993.
	
	\bibitem{Ekeland}
	I. Ekeland,
	{\em On the variational principle},
	J. Math. Anal. Appl. {\bf 47}, (1974), 324--353.
	
	\bibitem{FelliUgu}
	V. Felli, F. Uguzzoni, 
	{\em Some existence results for the Webster scalar curvature problem in presence of symmetry},
	Ann. Mat. Pura Appl. {\bf 183}(4), (2004), 469--493.
	
%	\bibitem{FlynnVetois}
%	J. Flynn, J. V\'{e}tois, 
%	{\em Liouville-type result for the CR Yamabe equation in the Heisenberg group},
%	to appear Ann. Sc. Norm. Super. Pisa Cl. Sci. (5), (2024), \url{https://doi.org/10.2422/2036-2145.202311_016}
	
	\bibitem{Folland}
	G.B. Folland,
	{\em Subelliptic estimates and function spaces on nilpotent Lie groups}, 
	Ark. Mat. {\bf 13}, (1975), 161--207.
	
%	\bibitem{Folland_per_Stein}
%	G.B. Folland, 
%		{\em The Heisenberg group and its relatives in the work of Elias M. Stein},
%	J. Geom. Anal. {\bf 31}(7), (2021), 6681--6697.
	
%	\bibitem{FollStein_CPAM}
%	G.B. Folland, E.M. Stein, 
%	{\em Estimates for the $\overline{\partial}_{b}$ complex and analysis on the Heisenberg group}, Commun. Pure Appl. Math. {\bf 27}, (1974), 429-–522. (1974) 
	
	\bibitem{FollandStein}
	G.B. Folland, E.M. Stein,
	{\em Hardy {{Spaces}} on {{Homogeneous
				Groups}}}, Princeton University Press, Dec. 2020.
	
%	\bibitem{FSSC}
%	B. Franchi, R. Serapioni, F. Serra Cassano, 
%	{\em Approximation and imbedding theorems for weighted Sobolev spaces
%		associated with Lipschitz continuous vector fields}, 
%	Boll. Un. Mat. Ital. B {\bf 11}(7), (1997), 83--117.		
	
	%	\bibitem{Fulks}
	%	W. Fulks, J.S. Maybee, 
	%	{\em A singular non-linear equation},
	%	Osaka Math. J. {\bf 12},(1960), 1--19.
	
%	\bibitem{Gamara}
%	N. Gamara, 
%	{\em The CR Yamabe conjecture the case $n = 1$}, 
%	J. Eur. Math. Soc. {\bf 3}(2), (2001), 105--137.
%	
%	\bibitem{GaYa}
%	N. Gamara, R. Yacoub, 
%	{\em CR Yamabe conjecture the conformally flat case}, 
%	Pacific J. Math. {\bf 201}(1), (2001), 121--175.
%	
%	\bibitem{GaUg}
%	E. Garagnani, F. Uguzzoni, 
%	{\em A multiplicity result for a degenerate elliptic equation with critical growth on noncontractible domains}, 
%	Topol. Methods Nonlinear Anal. {\bf 22}(1), (2003), 53--68.
	
	%	\bibitem{Garain}
	%P. Garain,
	%{\em On a class of mixed local and nonlocal semilinear elliptic equation with singular nonlinearity}, J. Geom. Anal. {\bf 33}, 212, (2023). 
	
	\bibitem{GaLa}
	N. Garofalo, E. Lanconelli, 
	{\em Existence and nonexistence results for
		semilinear equations on the Heisenberg group},
	Indiana Univ. Math. J. {\bf 41}(1), (1992), 71--98.

\bibitem{GaroVa2}
	N. Garofalo, D. Vassilev, 
	{\em Regularity near the characteristic set in the non-linear Dirichlet problem and conformal geometry of sub-Laplacians on Carnot groups}, 
	Math Ann {\bf 318}, (2000), 453–-516.
	
	
	\bibitem{GaroVa}
	N. Garofalo, D. Vassilev, 
	{\em Symmetry properties of positive entire solutions of Yamabe-type equations on groups of Heisenberg type}, Duke Math. J. {\bf 106}(3), (2001), 411--448.
	
	
	%\bibitem{GilbargTrudinger}
	%D.~Gilbarg and N.~S. Trudinger,
	%{\em Elliptic {{Partial Differential
	%			Equations}} of {{Second Order}}}, vol.~224 of Classics in {{Mathematics}},
	%Springer, Berlin, Heidelberg, 2001.
	
	
	%\bibitem{Giusti}
	%  E. Giusti, 
	% \emph{Direct Methods in the Calculus of Variations}, 
	% World Scientific Publishing Co., Inc., River Edge, NJ, 2003.
	
%	\bibitem{GutLan}
%	C.E. Guti\'{e}rrez, E. Lanconelli,
%	{\em Maximum principle, nonhomogeneous Harnack inequality, and Liouville theorems for  X-elliptic operators},	
%	Comm. Partial Differential Equations {\bf 28}(11-12), (2003), 1833--1862.
%	
%	\bibitem{Haitao}
%	Y. Haitao,
%	{\em Multiplicity and asymptotic behavior of positive solutions for a singular semilinear elliptic problem},
%	J. Differential Equations {\bf 189}(2), (2003), 487--512.
%	
	
	\bibitem{Hormander}
	L. H\"{o}rmander, 
	{\em Hypoelliptic second order differential equations}, 
	Acta Math. {\bf 119}, (1967), 147--171.
	
	\bibitem{Jerison}
	D.S. Jerison, 
	{\em The Dirichlet problem for the Kohn Laplacian on the Heisenberg group. I}, 
	J. Functional Analysis {\bf 43}(1), (1981), 97--142.
	
	\bibitem{Jerison2}
	D.S. Jerison, 
	{\em The Dirichlet problem for the Kohn Laplacian on the Heisenberg group. II}, 
	J. Functional Analysis {\bf 43}(2), (1981), 224--257.
	
	
	\bibitem{JerisonLee}
	D.S. Jerison, J.M. Lee, 
	{\em The Yamabe problem on CR manifolds}, 
	J. Differential Geom. {\bf 25}(2), (1987), 167--197.
	
	\bibitem{JerisonLee2}
	D.S. Jerison, J.M. Lee, 
	{\em Extremals for the Sobolev inequality on the Heisenberg group and the CR Yamabe problem}, 
	J. Amer. Math. Soc. {\bf 1}(1), (1988), 1--13.

	\bibitem{JerisonLee3}
D.S. Jerison, J.M. Lee, 
{\em Intrinsic CR normal coordinates and the CR Yamabe problem},
J. Differ. Geom. {\bf 29}, (1989), 303--343.
%	
%	\bibitem{KS}
%	D. Kinderlehrer, G. Stampacchia, 
%	{\em An introduction to variational inequalities and their applications}, 
%	Volume 88 of Pure and Applied Mathematics. Academic Press,
%	Inc. Harcourt Brace Jovanovich, Publishers, New York-London, 1980.	
	
%	\bibitem{Kohn}	
%	J.J. Kohn, 
%	{\em Boundaries of complex manifolds}, 
%	Proc. Conf. Complex Analysis (Minneapolis, 1964),
%	Springer, Berlin-Heidelberg-New York, (1965), 81--94.
	
%	\bibitem{MaOu}
%	X.-N. Ma, Q. Ou,
%	{\em A Liouville theorem for a class semilinear elliptic equations on the Heisenberg group}. 
%	Adv. Math. {\bf 413}(15), (2023), 1--20.
	

	
%	\bibitem{LaUg}
%	E. Lanconelli, F. Uguzzoni, 
%	{\em Non-existence results for semilinear Kohn–Laplace equations in unbounded domains}, 
%	Comm. Partial Differential Equations {\bf 25}, (2000), 1703--1739.
%	
% 
	
	%	\bibitem{LeoniFract}
	%	G. Leoni,
	%	\emph{A First Course in Fractional Sobolev Spaces}, 
	%	Graduate Studies in Mathematics (2023).
	
	%	\bibitem{Lions}
	%	P.L. Lions, 
	%	{\em The concentration-compactness principle in the calculus of variations. The limite case. II}, 
	%	Rev. Mat. Iberoam {\bf 1}, (1985), 45--121.
	
	\bibitem{Loiudice1}
	A. Loiudice,
	{\em Semilinear subelliptic problems with critical growth
		on Carnot groups}, Manuscripta Math. {\bf 124}, (2007), 247--259.
	
	%		\bibitem{Loiudice2}
	%	A. Loiudice,
	%		{\em $L^p$-weak regularity and asymptotic behavior of solutions for critical equations with singular potentials on Carnot groups},
	%	NoDEA Nonlinear Differential Equations Appl. {\bf 17}(5), (2010), 575--589.
	
	\bibitem{Loiudice2}
	A. Loiudice,
	{\em Critical growth problems with singular nonlinearities on Carnot groups},
	Nonlinear Anal. {\bf 126}, (2015), 415--436.
	
	\bibitem{Loiudice3}
	A. Loiudice,
	{\em Optimal decay of $p$-Sobolev extremals on Carnot groups},
	J. Math. Anal. Appl. {\bf 470}(1), (2019), 619--631.
	
	\bibitem{Loiudice4}
	A. Loiudice,
	{\em Critical problems with Hardy potential on stratified Lie groups},
	Adv. Differential Equations {\bf 28}(1-2), (2023), 1--33.
	
%	\bibitem{LuWei}
%	G. Lu, J. Wei, 
%	{\em On positive entire solutions to the Yamabe-type problem on the Heisenberg and stratified
%		groups}, 
%	Electron. Res. Announc. Amer. Math. Soc. {\bf 3}, (1997), 83--89.
	
	
%	\bibitem{MaMa}
%	A. Maalaoui, V. Martino,
%	{\em Multiplicity result for a nonhomogeneous Yamabe type equation involving the Kohn Laplacian},
%	J. Math. Anal. Appl. {\bf 399}(1), (2013), 333--339.
	
	\bibitem{MaMaPi}
	A. Maalaoui, V. Martino, A. Pistoia, 
	{\em Concentrating solutions for a sub-critical sub-elliptic problem},
	Differential Integral Equations {\bf 26}(11-12), (2013), 1263--1274.	
%	
%	\bibitem{MaMaTr}
%	A. Maalaoui, V. Martino, G. Tralli,
%	{\em Complex group actions on the sphere and sign changing solutions for the CR-Yamabe equation},
%	J. Math. Anal. Appl. {\bf 431}(1), (2015), 126--135.
	
	\bibitem{MaUg}
	A. Malchiodi, F. Uguzzoni, 
	{\em A perturbation result for the Webster scalar curvature problem on the CR sphere},
	J. Math. Pures Appl. {\bf 81}(10), (2002), 983--997.
	
	%		
	%\bibitem{MolicaFerrara}		
	%G. Molica Bisci, M. Ferrara, 
	%{\em Subelliptic and parametric equations on Carnot groups}, Proc.
	%Amer. Math. Soc.{\bf 144}, (2016), 3035--3045.
	
	%	\bibitem{MolicaPucci}
	%	G. Molica Bisci, P. Pucci, 
	%	{\em Critical Dirichlet problems on H domains of Carnot groups}, 
	%	Electron. J.
	%	Diff. Equ. {\bf 25}, (2017), 179--196.	
	%	
	
%	\bibitem{MolicaRepovs}
%	G. Molica Bisci, D. Repovs, 
%	{\em Yamabe-type equations on Carnot groups}, 
%	Potential Anal. {\bf 46}(2), (2017), 369--383.
	
	%\bibitem{NaStWa}
	%A. Nagel, E.M. Stein, S. Wainger, 
	%{\em Balls and metrics defined by vector fields I: basic properties}, 
	%Acta Math. {\bf 155} (1985), 103--147. 
	

	
	\bibitem{PaPiTe}
	G. Palatucci, M. Piccinini, L. Temperini,	
	{\em Struwe's global compactness and energy approximation of the critical Sobolev embedding in the Heisenberg group},
	Adv. Calc. Var. (2024), \url{https://doi.org/10.1515/acv-2024-0044}
	
%	\bibitem{Prajapat}
%	J.V. Prajapat, A.S. Varghese,
%	{\em Symmetry and classification of solutions to an integral equation in the Heisenberg group},
%	Math. Ann. {\bf 392}, (2025), 659--700.
	
%	\bibitem{RotSte}
%	L.P. Rothschild, E.M. Stein, 
%	{\em Hypoelliptic differential operators and nilpotent groups},
%	Acta Math. {\bf 137}(3-4), (1976), 247--320.
	
	\bibitem{Ruzhansky_Suragan}
    M. Ruzhansky, D. Suragan, 
    {\em Green’s Identities, Comparison Principle and Uniqueness of Positive Solutions for Nonlinear p-sub-Laplacian Equations on Stratified Lie Groups}, 
    Potential Anal {\bf 53}, (2020), 645--658.
	
%	\bibitem{Semmes}
%	S. Semmes, 
%	{\em On the nonexistence of bi-Lipschitz parameterizations and geometric problems about
%		A-weights}, 
%	Rev. Mat. Iberoam. {\bf 12}(2), (1996), 337--410.
	
	\bibitem{Struwe}
	M. Struwe, 
	{\em Variational Methods}, 
	Springer-Verlag, Berlin, 1990. xiv+244 pp.
	

	
	
	\bibitem{Tarantello}
	G. Tarantello, 
	{\em On nonhomogeneous elliptic equations involving critical Sobolev exponent}, 
	Ann. Inst. H. Poincar\'{e} C Anal. Non Lin\'{e}aire {\bf 9}, (1992), 281--304.
	
	\bibitem{Ugu1}
	F. Uguzzoni, 
	{\em A non-existence theorem for a semilinear Dirichlet problem involving critical exponent on halfspaces of the Heisenberg group}, 
	NoDEA Nonlinear Differential Equations Appl. {\bf 6}, (1999), 191--206.
	
%	\bibitem{Ugu2}
%	F. Uguzzoni, 
%	{\em A note on Yamabe-type equations on the Heisenberg group}, Hiroshima Math. J. {\bf 30}, (2000), 179--189.
	
	%	\bibitem{Wang}
	%	W. Wang,
	%	{\em Positive Solution of a Subelliptic Nonlinear
		%		Equation on the Heisenberg Group},
	%	Canad. Math. Bull. Vol. {\bf 44}(3), (2001), 346--354.
	
	
	
\end{thebibliography}


 \end{document}
